%%%%%%%%%%%%%%%%%%%%%%%%%%%%%%%%%%%%%%%%%%%%%%%%%%%%%%%%%%%%%%%%%%%%%
%                                                                   %
%                      Chapter 1 - Section 1                        %
%                                                                   %
%%%%%%%%%%%%%%%%%%%%%%%%%%%%%%%%%%%%%%%%%%%%%%%%%%%%%%%%%%%%%%%%%%%%% 

\chapter{A Context for Calculus}

Calculus gives us a language to describe how quantities are related to
one another, and it gives us a set of computational and visual tools
for exploring those relationships.  Usually, we want to understand how
quantities are related in the context of a particular problem---it
might be in chemistry, or public policy, or mathematics itself.  In
this chapter we take a single context---an infectious disease
spreading through a population---to see how calculus emerges and how
it is used.


\section{The Spread of Disease}

\subsection{Making a Model}

Many human diseases are contagious: you ``catch'' them from someone
who is already infected.  Contagious diseases are of many kinds.
Smallpox, polio, and plague are severe and even fatal, 
%
\mar{Some properties of contagious diseases} 
%
while the common cold and the childhood illnesses of measles, mumps,
and rubella are usually relatively mild.  Moreover, you can catch a
cold over and over again, but you get measles only once.  A disease
like measles is said to ``confer immunity'' on someone who recovers
from it.  Some diseases have the potential to affect large segments of
a population; they are called {\em epidemics}\index{epidemic} (from
the Greek words {\em epi}, upon + {\em demos}, the people.) {\em
  Epidemiology\/}\index{epidemiology} is the scientific study of these
diseases.

An epidemic is a complicated matter, but the dangers posed by
contagion---and especially by the appearance of new and uncontrollable
diseases---compel us to learn as much as we can about the nature of
epidemics.  Mathematics offers a very special kind of help.  First, we
can try to draw out of the situation its essential features and
describe them mathematically.  This is calculus as {\em language}.  We
substitute an ``ideal'' mathematical world for the real one.
%
\mar{The idea of a mathematical model}
%

This is calculus as {\em tool}.  Any conclusion we reach about the
model can then be interpreted to tell us something about the reality.

To give you an idea how this process works, we'll build a model of an
epidemic.  Its basic purpose is to help us understand the way a
contagious disease spreads through a population---to the point where
we can even predict what fraction falls ill, and when.  Let's suppose
the disease we want to model is like measles.  In particular,
\begin{itemize}
 \item it is mild, so anyone who falls ill eventually recovers;
 \item it confers permanent immunity on every recovered victim.
\end{itemize}
In addition, we will assume that the affected population is large but
fixed in size and confined to a geographically well-defined region.
To have a concrete image, you can imagine the elementary school
population of a big city.

At any time, that population can be divided into three distinct
classes:
\begin{description}
 \item[Susceptible:]\index{susceptible} those who have never had the
   illness and can catch it;
 \item[Infected:]\index{infected} those who currently have the illness
   and are contagious;
 \item[Recovered:]\index{recovered} those who have already had the
   illness and are immune.
\end{description}

Suppose we let $S$, $I$, and $R$ 
%
\mar{The quantities that our model analyzes}
%
denote the number of people in each of these three classes,
respectively.  Of course, the classes are all mixed together
throughout the population: on a given day, we may find persons who are
susceptible, infected, and recovered in the same family.  For the
purpose of organizing our thinking, though, we'll represent the whole
population as separated into three ``compartments'' as in the
following diagram:\index{compartment diagram}

\setlength{\unitlength}{1.5pt}
\begin{picture}(180,110)(-109,-55)
	\bezier{150}(-30,18)(0,68)(30,18) %150
	\bezier{150}(30,18)(62,-35)(0,-35) %150
	\bezier{150}(0,-35)(-62,-35)(-30,18) %150
	\put(0,0){\line(5,3){30}}
	\put(0,0){\line(-5,3){30}}
	\put(0,0){\line(0,-1){35}}
	\put(0,20){\makebox[0pt]{$S$}}
	\put(20,-12){\makebox[0pt]{$I$}}
	\put(-20,-12){\makebox[0pt]{$R$}}
	\put(0,48){\makebox[0pt]{Susceptible}}
	\put(46,-20){Infected}
	\put(-80,-20){Recovered}
	\put(0,-45){\makebox[0pt]{The Whole Population}}
\end{picture}

The goal of our model is to determine what happens to the numbers $S$,
$I$, and $R$ over the course of time.  Let's first see what our
knowledge and experience of childhood diseases might lead us to
expect.  When we say there is a ``measles outbreak,'' we mean that
there is a relatively sudden increase in the number of cases, and then
a gradual decline.  After some weeks or months, the illness virtually
disappears.  In other words, the number $I$ is a {\bf
  variable};\index{variable} its value changes over time.  One
possible way that $I$ might vary is shown in the following
graph.\index{graph}\label{Sgraph}


\setlength{\unitlength}{1pt}
\begin{picture}(340,140)(15,0)
	\put(50,20){\vector(1,0){255}}
	\put(50,20){\vector(0,1){100}}
	\put(290,8){time}
	\put(39,114){$I$}
	\begin{thicklines}
		\bezier{75}(50,25)(65,35)(75,60)
		\bezier{100}(75,60)(99,120)(124,72) %100
		\bezier{200}(124,72)(149,24)(300,23) %200
	\end{thicklines}
\end{picture}

During the course of the epidemic, susceptibles are constantly falling
ill.  Thus we would expect the number $S$ to show a steady decline.
Unless we know more about the illness, we cannot decide whether
everyone eventually catches it.  In graphical terms, this means we
don't know whether the graph of $S$ levels off at zero or at a value
above zero.  Finally, we would expect more and more people in the
recovered group as time passes.  The graph of $R$ should therefore
climb from left to right.  The graphs of $S$ and $R$ might take the
following forms:


\begin{picture}(340,140)(15,0)
	\put(50,20){\vector(1,0){255}}
	\put(50,20){\vector(0,1){100}}
	\put(290,8){time}
	\put(38,114){$S$}
	\begin{thicklines}
	\bezier{80}(50,100)(70,94)(100,71)
	\bezier{200}(100,71)(160,25)(300,24) %200
	\end{thicklines}
\end{picture}


\begin{picture}(340,140)(15,0)
	\put(50,20){\vector(1,0){255}}
	\put(50,20){\vector(0,1){100}}
	\put(290,8){time}
	\put(37,114){$R$}
	\begin{thicklines}
	\bezier{80}(50,30)(90,40)(120,70)
	\bezier{200}(120,70)(160,110)(300,110) %200
	\end{thicklines}
\end{picture}
\label{Rgraph}

While these graphs give us an idea about what might happen, they raise
some new questions, too.  For example, because there are no scales
marked along the axes,
%
\mar{Some quantitative questions that the graphs raise}
%
the first graph does not tell us how large $I$ becomes when the
infection reaches its peak, nor when that peak occurs.  Likewise, the
second and third graphs do not say how rapidly the population either
falls ill or recovers.  A good model of the epidemic should give us
graphs like these and it should also answer the quantitative questions
we have already raised---for example: When does the infection hit its
peak?  How many susceptibles eventually fall ill?


\subsection{A Simple Model for Predicting Change}

Suppose\label{`modelstart'} we know the values of $S$, $I$, and $R$
today; can we figure out what they will be tomorrow, or the next day,
or a week or a month from now?  Basically, this is a problem of
predicting the future.  One way to deal with it is to get an idea how
$S$, $I$, and $R$ are {\em changing}.  To start with a very simple
example, suppose the city's Board of Health reports that the measles
infection has been spreading at the rate of 470 new cases per day for
the last several days.  If that rate continues to hold steady and we
start with 20,000 susceptible children, then we can expect 470 fewer
susceptibles with each passing day.  The immediate future would then
look like this:
\begin{center}
\begin{tabular}{ccc}
			& accumulated		& remaining \\ [-2pt]
	days after& number of new	& number of \\ [-2pt]
	today	& infections		& susceptibles \\[2pt] 
		\hline 
	\begin{tabular}{c} \rule{0pt}{15pt} 0 \rule{0pt}{15pt} 
		\\ 1 \\ 2 \\ 3 \end{tabular} &
	\begin{tabular}{r} \rule{0pt}{15pt} 0 \\ 470 \\ 940 \\
		$1410$ \end{tabular}	&
	\begin{tabular}{r} \rule{0pt}{15pt}$20\,000$\\$19\,530$
	 \\ 	$19\,060$ \\ $18\,590$ \end{tabular} \\[-3pt]
	\vdots	&	\vdots		&	\vdots
\end{tabular}
\end{center}

Of course, these numbers will be correct only if the infection
continues to spread at its present rate of 470 persons per day.
%
\mar{Knowing rates, we can predict future values}  
%
If we want to follow $S$, $I$, and $R$ into the future, our example
suggests that we should pay attention to the {\bf
  rates}\index{rate}\index{rate of change} at which these quantities
change.  To make it easier to refer to them, let's denote the three
rates by $S'$, $I'$, and $R'$.  For example, in the illustration
above, $S$ is changing at the rate $S' = -470$ persons per day.  We
use a minus sign here because $S$ is {\em decreasing\/} over time.  If
$S'$ stays fixed we can express the value of $S$ after $t$ days by the
following formula:
\[
S = 20000 + S' \cdot t = 20000 - 470\,t \mbox{ persons.}
\] 
Check that this gives the values of $S$ found in the table when $t =
0$, 1, 2, or 3.  How many susceptibles does it say are left after 10
days?

Our assumption that $S' = -470$ persons per day amounts to a
mathematical characterization
%
\mar{The equation $S' = -470$ \\ is a model}
%
of the susceptible population---in other words, a model!  Of course it
is quite simple, but it led to a formula that told us what value we
could expect $S$ to have at any time $t$.

The model will even take us backwards in time\label{backwards}.  For
example, two days ago the value of $t$ was $-2$; according to the
model, there were
\[
S = 20000 - 470 \times -2 = 20940
\]
susceptible children then.  There is an obvious difference between
going backwards in time and going forwards: we already know the past.
Therefore, by letting $t$ be negative we can generate values for $S$
that can be checked against health records.  If the model gives good
agreement with {\em known\/} values of $S$ we become more confident in
using it to predict {\em future\/} values.

To predict the value of $S$ using the rate $S'$ we clearly need to
have a starting point---a known value of $S$
%
\mar{Predictions depend on the initial value, too}
%
from which we can measure changes.  In our case that starting point is
$S = 20000$.  This is called the {\bf initial value}\index{initial value} 
of $S$, because it is given to us at the ``initial time'' $t
= 0$.  To construct the formula $S = 20000 - 470\,t$, we needed to
have an initial value as well as a rate of change for $S$.

\smallskip

In the following pages we will develop a more complex model for all
three population groups that has the same general design as this
simple one.  Specifically, the model will give us information about
the rates $S'$, $I'$, and $R'$, and with that information we will be
able to predict the values of $S$, $I$, and $R$ at any time
$t$\label{`modelend'}.


\subsection{The Rate of Recovery}
\index{rate!of recovery}
\index{recovery}

Our first task will be to model the recovery rate $R'$.  We look at
the process of recovering first, because it's simpler to analyze.  An
individual caught in the epidemic first falls ill and then
recovers---recovery is just a matter of time.  In particular, someone
who catches measles\index{measles} has the infection for about
fourteen days.  So if we look at the entire infected population today,
we can expect to find some who have been infected less than one day,
some who have been infected between one and two days, and so on, up to
fourteen days.  Those in the last group will recover today.  In the
absence of any definite information about the fourteen groups, let's
assume they are the same size.  Then 1/14-th of the infected
population will recover today:
\[
\mbox{today the change in the recovered population} = 
          \dfrac{I \mbox{ persons}}{14 \mbox{ days}}.
\]
There is nothing special about today, though; $I$ has a value at any
time.  Thus we can make the same argument about any other day:
\[
\mbox{every day the change in the recovered population} 
		= \dfrac{I \mbox{ persons}}{14 \mbox{ days}}.
\]
This equation is telling us about $R'$, the rate at which $R$ is
changing.  We can write it more simply in the form
%
\mar{\vspace{.2in} The first piece of the $S$-$I$-$R$ model}
%
\[
R' = \dfrac{1}{14}I \mbox{ persons per day} .
\]

\newpage

We call this a {\bf rate equation}.\index{rate equation} Like any
equation, it links different quantities together.  In this case, it
links $R'$ to $I$.  The rate equation for $R$ is the first part of our
model of the measles epidemic.

\aside{Are you uneasy about our claim that {1/14-th} of the infected
  population recovers every day?  You have good reason to be.  After
  all, during the first few days of the epidemic almost no one has had
  measles the full fourteen days, so the recovery rate will be much
  less than $I/14$ persons per day.  About a week before the infection
  disappears altogether there will be no one in the early stages of
  the illness.  The recovery rate will then be much greater than
  $I/14$ persons per day.  Evidently our model is not a perfect mirror
  of reality!

\quad \ Don't be particularly surprised or dismayed by this.  Our goal
is to gain insight into the workings of an epidemic and to suggest how
we might intervene to reduce its effects.  So we start off with a
model which, while imperfect, still captures some of the workings.
The simplifications in the model will be justified if we are led to
inferences which help us understand how an epidemic works and how we
can deal with it.  If we wish, we can then refine the model, replacing
the simple expressions with others that mirror the reality more
fully.}

Notice that the rate equation for $R'$ does indeed give us a tool to
predict future values of $R$.  For suppose today 2100 people are
infected and 2500 have already recovered.  Can we say how large the
recovered population will be tomorrow or the next day?  Since $I =
2100$,
\[
R' = \dfrac{1}{14} \times 2100 = 150 \mbox{ persons per day}\, .
\]
Thus 150 people will recover in a single day, and twice as many, or
300, will recover in two.  At this rate the recovered population will
number 2650 tomorrow and 2800\label{`future'} the next day.

These calculations assume that the rate $R'$ holds steady at 150
persons per day for the entire two days.  Since $R' = I/14$, this is
the same as assuming that $I$ holds steady at 2100 persons.  If
instead $I$ varies during the two days we would have to adjust the
value of $R'$ and, ultimately, the future values of $R$ as well.  In
fact, $I$ {\em does\/} vary over time.  We shall see this when we
analyze how the infection is transmitted.  Then, in chapter 2, we'll
see how to make the adjustments in the values of $R'$ that will permit
us to predict the value of $R$ in the model with as much accuracy as
we wish.

\medskip

\noindent {\bf Other diseases}.  What can we say about the recovery
rate for a contagious disease other than measles? If the period of
infection of the new illness is $k$ days, instead of 14, and if we
assume that 1/$k$ of the infected people recover each day, then the
new recovery rate is
\[
R' = \dfrac{I \mbox{ persons}}{k \mbox{ days}} = 
			\dfrac{1}{k}I  \mbox{ persons per day.}
\]
If we set $b = 1/k$ we can express the recovery rate equation in the
form
\[
R' = bI  \mbox{ persons per day.}
\]
The constant $b$ is called the {\bf recovery
  coefficient}\index{recovery!coefficient} in this context.

\smallskip
\noindent
\begin{minipage}[b]{1.1in}
%\setlength{\unitlength}{1.5pt}
\begin{picture}(110,70)(-12,-30)
	\bezier{200}(-30,39)(0,89)(30,39) %200
	\put(0,21){\line(5,3){30}}
	\put(0,21){\line(-5,3){30}}
	\put(0,41){\makebox[0pt]{\footnotesize $S$}}

	\bezier{200}(42,18)(74,-35)(12,-35) %200
	\put(12,0){\line(5,3){30}}
	\put(12,0){\line(0,-1){35}}
	\put(32,-12){\makebox[0pt]{\footnotesize $I$}}

	\bezier{200}(-12,-35)(-74,-35)(-42,18) %200
	\put(-12,0){\line(-5,3){30}}
	\put(-12,0){\line(0,-1){35}}
	\put(-32,-12){\makebox[0pt]{\footnotesize $R$}}

%	\put(-11,-18){\line(1,1){9}}
%	\put(-11,-18){\line(1,-1){9}}
%	\multiput(-11,-18)(.3,.3){30}{\circle*{.1}}
%	\multiput(-11,-18)(.3,-.3){30}{\circle*{.2}}
	\bezier{30}(-11,-18)(-6.5,-13.5)(-2,-9)
	\bezier{30}(-11,-18)(-6.5,-22.5)(-2,-27)
	\put(-2,-11){\line(0,1){2}}
	\put(-2,-11){\line(1,0){12}}
	\put(-2,-25){\line(0,-1){2}}
	\put(-2,-25){\line(1,0){12}}
	\put(10,-11){\line(0,-1){14}}
	\put(0,-21){\makebox[0pt]{\footnotesize $b$}}
	\put(0,-45){\makebox[0pt]{\footnotesize recovery}}
\end{picture}
\setlength{\unitlength}{1pt}
\end{minipage} \hfill
\begin{minipage}[b]{3.8in}
\quad\ Let's incorporate our understanding of recovery into the
compartment diagram.  For the sake of illustration, we'll separate the
three compartments.  As time passes, people ``flow'' from the infected
compartment to the recovered.  We represent this flow by an arrow from
$I$ to $R$.  We label the arrow with the recovery coefficient $b$ to
indicate that the flow is governed by the rate equation $R' = b I$.
\end{minipage}


\subsection{The Rate of Transmission}
\index{rate!of transmission}
\index{transmission}

Since susceptibles become infected, the compartment diagram above
should also have an arrow that goes from $S$ to $I$ and a rate
equation for $S'$ to show how $S$ changes as the infection spreads.
While $R'$ depends only on $I$, because recovery involves only waiting
for people to leave the infected population, $S'$ will depend on both
$S$ and $I$, because transmission involves contact between susceptible
and infected persons.

Here's a way to model the transmission rate.  First, consider a single
susceptible person on a single day.  On average, this person will
contact only a small fraction, $p$, of the infected population.  For
example, suppose there are 5000 infected children, so $I = 5000$.  We
might expect only a couple of them---let's say 2---will be in the same
classroom with our ``average'' susceptible.  So the fraction of
contacts is $p = 2/I = 2/5000 = .0004$.  The 2 contacts themselves can
be expressed as $2 = (2/I)\cdot I = pI$ contacts per day per
susceptible.

To find out how many daily contacts the {\em whole\/} 
%
\mar{Contacts are proportional to both $S$ and $I$}
%
susceptible population will have, we can just multiply the average
number of contacts per susceptible person by the number of
susceptibles: this is $pI \cdot S = pSI$.

Not all contacts lead to new infections; only a certain fraction $q$
do.  The more contagious the disease, the larger $q$ is.  Since the
number of daily contacts is $pSI$, we can expect $q \cdot pSI$ new
infections per day (i.e., to convert contacts to infections, multiply
by $q$).  This becomes $aSI$ if we define $a$ to be the product $qp$.

Recall, the value of the recovery coefficient $b$ depends only on the
illness involved.  It is the same for all populations.  By contrast,
the value of $a$ depends on the general health of a population and the
level of social interaction between its members.  Thus, when two
different populations experience the same illness, the values of $a$
could be different.  One strategy for dealing with an epidemic is to
alter the value of $a$.  Quarantine\index{quarantine} does this, for
instance; see the exercises.

Since each new infection decreases the number of susceptibles, we have
the rate equation for $S$: 
%
\mar{\vspace{4pt} Here is the second piece of the $S$-$I$-$R$ model}
%
\[
S' =  - a S I \mbox{ persons per day} \, .
\]
The minus sign here tells us that $S$ is decreasing (since $S$ and $I$
are positive).  We call $a$ the {\bf transmission
  coefficient}.\index{transmission!coefficient}

\smallskip
\noindent
\begin{minipage}[b]{3.8in}
\quad\ Just as people flow from the infected to the recovered
compartment when they recover, they flow from the susceptible to the
infected when they fall ill.  To indicate the second flow let's add
another arrow to the compartment diagram.
\index{compartment diagram}\label{`diagram'} Because this flow is due to the
transmission of the illness, we will label the arrow with the
transmission coefficient $a$.  The compartment diagram now reflects
all aspects of our model.
\end{minipage}
\begin{minipage}[b]{1.1in}
\begin{picture}(110,70)(-74,-30)
	\bezier{200}(-30,39)(0,89)(30,39) %200
	\put(0,21){\line(5,3){30}}
	\put(0,21){\line(-5,3){30}}
	\put(0,41){\makebox[0pt]{\footnotesize $S$}}

	\put(27,10){\line(-4,1){12}}
	\put(27,10){\line(1,4){3}}
	\put(10,24){\line(3,-5){6}}
	\put(10,24){\line(5,3){12}}
	\put(22,31.2){\line(3,-5){6}}
	\put(21,18){\makebox[0pt]{\footnotesize $a$}}
	\put(37,25){\footnotesize transmission}

	\bezier{200}(42,18)(74,-35)(12,-35) %200
	\put(12,0){\line(5,3){30}}
	\put(12,0){\line(0,-1){35}}
	\put(32,-12){\makebox[0pt]{\footnotesize $I$}}

	\bezier{200}(-12,-35)(-74,-35)(-42,18) %200
	\put(-12,0){\line(-5,3){30}}
	\put(-12,0){\line(0,-1){35}}
	\put(-32,-12){\makebox[0pt]{\footnotesize $R$}}

	\bezier{30}(-11,-18)(-6.5,-13.5)(-2,-9)
	\bezier{30}(-11,-18)(-6.5,-22.5)(-2,-27)
	\put(-2,-11){\line(0,1){2}}
	\put(-2,-11){\line(1,0){12}}
	\put(-2,-25){\line(0,-1){2}}
	\put(-2,-25){\line(1,0){12}}
	\put(10,-11){\line(0,-1){14}}
	\put(0,-21){\makebox[0pt]{\footnotesize $b$}}
	\put(0,-45){\makebox[0pt]{\footnotesize recovery}}
\end{picture}
\setlength{\unitlength}{1pt}
\end{minipage}

\aside{We haven't talked about the units in which to measure $a$ and
  $b$.  They must be chosen so that any equation in which $a$ or $b$
  appears will balance.  Thus, in $R' = bI$ the units on the left are
  persons/day; since the units for $I$ are persons, the units for $b$
  must be 1/(days).  The units in $S' = -aSI$ will balance only if $a$
  is measured in 1/(person-day).

\quad\ The reciprocals have more natural interpretations.  First of
all, 1/$b$ is the number of days a person needs to recover.  Next,
note that 1/$a$ is measured in person-days (i.e., persons $\times$
days), which are the natural units in which to measure exposure.  Here
is why.  Suppose you contact 3~infected persons for each of 4~days.
That gives you the same exposure to the illness that you get from
6~infected persons in 2~days---both give 12 ``person-days'' of
exposure.  Thus, we can interpret $1/a$ as the level of exposure of a
typical susceptible person.}


\subsection{Completing the Model}

The final rate equation we need---the one for $I'$---reflects what is
already clear from the compartment diagram: every loss in $I$ is due
to a gain in $R$, while every gain in $I$ is due to a loss in $S$.

\newpage

\mbox{}\vspace{-12pt}
%
\mar{\vspace{10pt} Here is the complete $S$-$I$-$R$ model}
%
\begin{alignat*}{2}
		S' &=  -a S I,           &    &     \\
		I' &=  \hphantom{-}a S I &-\, & b I,  \\
		R' &=                    &    & b I.
\end{alignat*}
If you add up these three rates you should get the overall rate of
change of the whole population.  The sum is zero.  Do you see why?

\aside{You should not draw the conclusion that the only use of rate
  equations\index{rate equation} is to model an epidemic.  Rate
  equations have a long history, and they have been put to many uses.
  Isaac Newton\index{Newton, I.} (1642--1727) introduced them to model
  the motion of a planet around the sun.  He showed that the same rate
  equations could model the motion of the moon around the earth and
  the motion of an object falling to the ground.  Newton created
  calculus as a tool to analyze these equations.  He did the work
  while he was still an undergraduate---on an extended vacation,
  curiously enough, because a plague epidemic was raging at Cambridge!

\quad\ Today we use Newton's rate equations to control the motion of
earth satellites and the spacecraft that have visited the moon and the
planets.  We use other rate equations to model radioactive decay,
chemical reactions, the growth and decline of populations, the flow of
electricity in a circuit, the change in air pressure with
altitude---just to give a few examples.  You will have an opportunity
in the following chapters to see how they arise in many different
contexts, and how they can be analyzed using the tools of calculus.}

The following diagram summarizes, in a schematic way, the relation
between our model\index{model} and the reality it seeks to portray.

\begin{picture}(340,90)(10,-15)
	\put(-10,0)
	{
	\begin{picture}(120,70)
		\put(0,0){\framebox(120,50)
			{\shortstack{people fall ill\\[-2pt]
				and eventually\\[1pt] recover}}}
		\put(60,57){\makebox[0pt]{\bf Reality}}
	\end{picture}
	}
	\put(230,0)
	{
	\begin{picture}(120,70)
		\put(0,0){\framebox(120,50)
			{\shortstack{$S' = -aSI$ \\$I' = aSI - bI$ \\
				$R' = bI$}}}
		\put(60,57){\makebox[0pt]{\bf Mathematics}}
	\end{picture}
	}
	\put(125,32){\vector(1,0){90}}
	\put(215,18){\vector(-1,0){90}}
	\put(170,40){\makebox(0,0)[b]{\shortstack
		{Modelling \\ [-1pt] the reality}}}
	\put(170,13){\makebox(0,0)[t]{\shortstack
		{Interpreting \\ [-1pt] the mathematics}}}
\end{picture}


\noindent  
The diagram calls attention to several facts.  
%
\mar{The model is part of mathematics; it only approximates reality}
%
First, the model\index{model} is a part of mathematics.  It is
distinct from the reality being modelled.  Second, the model is based
on a simplified interpretation of the epidemic.  As such, it will not
match the reality exactly; it will be only an {\bf
  approximation}.\index{approximation} Thus, we cannot expect the
values of $S$, $I$, and $R$ that we calculate from the rate equations
to give us the exact sizes of the susceptible, infected, and recovered
populations.  Third, the connection between reality and mathematics is
a two-way street.  We have already travelled one way by constructing a
mathematical object that reflects some aspects of the epidemic.  This
is model-building.  Presently we will travel the other way.  First we
need to get mathematical answers to mathematical questions; then we
will see what those answers tell us about the epidemic.  This is
interpretation of the model.  Before we begin the interpretation, we
must do some mathematics.


\subsection{Analyzing the Model}

Now that we have a model\label{`analyzestart'} we shall analyze it as
a mathematical object.  We will set aside, at least for the moment,
the connection between the mathematics and the reality.  Thus, for
example, you should not be concerned when our calculations produce a
value for $S$ of 44,446.6 persons at a certain time---a value that
will never be attained in reality.  In the following analysis $S$ is
just a numerical quantity that varies with $t$, another numerical
quantity.  Using only mathematical tools we must now extract from the
rate equations the information that will tell us just how $S$ and $I$
and $R$ vary with $t$.

We already took the first steps in that direction when we used the
rate equation $R' = I/14$ to predict the value of $R$ two days into
the future (see page~\pageref{`future'}).  We assumed that $I$
remained fixed at 2100 during those two days, so the rate $R' =
2100/14 = 150$ was also fixed.  We concluded that if $R = 2500$ today,
it will be 2650 tomorrow and 2800 the next day.

A glance at the full $S$-$I$-$R$ model tells us those first steps have
to be modified.  The assumption
%
\mar{Rates are continually changing; this affects the calculations}
%
we made---that $I$ remains fixed---is not justified, because $I$ (like
$S$ and $R$) is continually changing.  As we shall see, $I$ actually
increases over those two days.  Hence, over the same two days, $R'$ is
not fixed at 150, but is continually increasing also.  That means that
$R'$ becomes larger than 150 during the first day, so $R$ will be
larger than $2650$ tomorrow.

The fact that the rates are continually changing complicates the
mathematical work we need to do to find $S$, $I$, and $R$.  In chapter
2 we will develop tools and concepts that will overcome this problem.
For the present we'll assume that the rates $S'$, $I'$, and $R'$ stay
fixed for the course of an entire day.  This will still allow us to
produce reasonable estimates\index{estimate} for the values of $S$,
$I$, and $R$.  With these estimates we will get our first glimpse of
the predictive power of the $S$-$I$-$R$ model.  We will also use the
estimates as the starting point for the work in chapter 2 that will
give us precise values.  Let's look at the details of a specific
problem.

\medskip

\noindent 
{\bf The Problem}.  Consider a measles epidemic in a school population
of 50,000 children.  The recovery coefficient is $b = 1/14$.  For the
transmission coefficient we choose $a = .00001$, a number within the
range used in epidemic studies.  We suppose that 2100 people are
currently infected and 2500 have already recovered.  Since the total
population is 50,000, there must be 45,400 susceptibles.  Here is a
summary of the problem in mathematical terms:
\begin{center}
\begin{picture}(360,140)
	\put(0,0){\framebox(360,140)
	{\begin{minipage}{4in}
	Rate equations:
	\vspace{-6pt}
	\begin{align*}
		S' &= -.00001SI,	\\
		I' &= .00001SI - I/14, \\
		R' &= I/14.	
		\end{align*}
	Initial values: when $t = 0$,
	\vspace{-6pt}
		$$
	S = 45400, \qquad I = 2100, \qquad R = 2500.
		$$
	\end{minipage}}}
\end{picture}
\end{center}

\smallskip

\noindent {\bf Tomorrow}.  From our earlier discussion, $R' = 2100/14
= 150$ persons per day, giving us an estimated value of $R = 2650$
persons for tomorrow.  To estimate $S$ we use
\[
S' = - .00001\,SI = - .00001 \times 45400 \times 2100
		= -953.4 \ \mbox{persons/day.}
\]
Hence we estimate that tomorrow 
\[
S = 45400 - 953.4 = 44446.6 \ \mbox{persons.}
\]
Since $S + I + R = 50000$ and we have $S + R = 47096.6$ tomorrow, a
final subtraction gives us $I = 2903.4$ persons.  (Alternatively, we
could have used the rate equation for $I'$ to estimate $I$.)

The fractional values in the estimates for $S$ and $I$ remind us that
the $S$-$I$-$R$ model describes the behavior of the epidemic only
approximately.

\medskip

\noindent {\bf Several days hence}.  According to the model, we
estimate that tomorrow $S = 44446.6$, $I = 2903.4$, and $R = 2650$.
Therefore, from the new $I$ we get a new approximation for the value
of $R'$ tomorrow; it is
\[
R' = \dfrac{1}{14} I = \dfrac{1}{14} \times 2903.4 = 207.4 \ 
			\mbox{persons/day.}
\]
Hence, two days from now we estimate that $R$ will have the new value
$2650 + 207.4 = 2857.4$.  Now follow this pattern to get new
approximations for $S'$ and $I'$, and then use those to estimate the
values of $S$ and $I$ two days from now.

\newpage

The pattern of steps that just carried you from the first day to the
second will work just as well to carry you from the second to the
third.
%
\mar{Stop and do\\ the calculations}
%
Pause now and do all these calculations yourself.  See exercises 15
and 16 on page~\pageref{`measles'}.  If you round your calculated
values of $S$, $I$, and $R$ to the nearest tenth, they should agree
with those in the following table.
\begin{center}
Estimates for the first three days\label{`threeday'}
\smallskip

{\small {\tabcolsep 3pt
\begin{tabular}{ccccccc}
	$t$ & $S$ & $I$ & $R$ & $S'$ & $I'$ & $R'$ \\ [2pt]
		\hline
	\begin{tabular}{c} \rule{0pt}{15pt} 0 \rule{0pt}{15pt} 
		\\ 1 \\ 2 \\ 3 \end{tabular} &
	\begin{tabular}{r} 
		\rule{0pt}{15pt} $45\,400.0$ \\ $44\,446.6$ \\
 		$43\,156.1$ \\ $41\,435.7$ \end{tabular} &
	\begin{tabular}{r} 
		\rule{0pt}{15pt} $2100.0$ \\ $2903.4$ \\
		 $3986.5$ \\ $5422.$1 \end{tabular} &
	\begin{tabular}{r} 
		\rule{0pt}{15pt} $2500.0$ \\ $2650.0$ \\
 		$2857.4$ \\ $3142.1$ \end{tabular} &
	\begin{tabular}{r} 
		\rule{0pt}{15pt} $-953.4$ \\ $-1290.5$ \\ 
		$-1720.4$ \\ \mbox{} \end{tabular} &
	\begin{tabular}{r} 
		\rule{0pt}{15pt} 803.4 \\ $1083.1$ \\ 
	$1435.7$ \\\mbox{} \end{tabular} &
	\begin{tabular}{r} 
		\rule{0pt}{15pt} 150.0 \\ 207.4 \\ 284.7 \\ 
		\mbox{} \end{tabular} 
\end{tabular}
}}
\end{center}

\noindent
{\bf Yesterday}.  We already pointed out, on
page~\pageref{backwards}, that we can use our models to go {\em
  backwards\/} in time, too.  This is a valuable way to see how well
the model fits reality, because we can compare estimates that the
model generates with health records for the days in the recent past.

To find how $S$, $I$, and $R$ change when we go one day into the
future we multiplied the rates $S'$, $I'$, and $R'$ by a time step of
+1.  To find how they change when we go one day into the past we do
the same thing, except that we must now use a time step of --1.
According to the table above, the rates at time $t = 0$ (i.e., today)
are
\[
S' = -953.4, \qquad I' = 803.4, \qquad R' = 150.0.
\]
Therefore we estimate that, one day ago, 
\begin{alignat*}{2}
S &= 45400 + (-953.4 \times -1)\; &= 45400 + 953.4 
		&=  46353.4,\\
I &= 2100 + (803.4 \times -1)   &= \hphantom{0}2100 - 803.4 
		&= 1296.6,  \\
R &= 2500 + (150.0 \times -1)   &= \hphantom{0}2500 - 150.0 
		&= 2350.0.  
\end{alignat*}
Just as we would expect with a spreading infection, there are more
susceptibles yesterday than today, but fewer infected or recovered.
In the exercises for section~3 you will have an opportunity to continue
this process, tracing the epidemic many days into the past.  For
example, you will be able to go back and see when the infection
started---that is, find the day when the value of $I$ was only about~1.

\newpage

\noindent
{\bf There and back again}.\index{there and back
  again}\label{`thereback'}
%
\mar{Go forward a day\\ and then back again}
%
What happens when we start with tomorrow's values and use tomorrow's
rates to go back one day---back to today?  We should get $S = 45400$,
$I = 2100$, and $R = 2500$ once again, shouldn't we?  Tomorrow's
values are
\begin{xalignat*}{3}
S  &= 44446.6, &   I  &= 2903.4, &   R  &= 2650.0, \\ 
S' &= -1290.5, &   I' &= 1083.1, &   R' &= 207.4.
\end{xalignat*}
To go backwards one day we must use a time step of $-1$.  The
predicted values are thus
\begin{alignat*}{2}
S &= 44446.6 + (-1290.6 \times -1) & &= 45737.2,  \\
I &= 2903.4 + (1083.1 \times -1) & &= 1820.3, \\
R &= 2650.0 + (207.4 \times -1) && = 2442.6. 
\end{alignat*}
These are {\em not\/} the values that we had at the start, when $t =
0$.  In fact, it's worth noting the difference\label{`differences'}
between the original values and those produced by ``going there and
back again.''
\begin{center}
\begin{tabular}{cccc}
		& original & there and	& \\[-2pt]
		& value	& back again	& difference \\ [2pt] \hline
	\begin{tabular}{c}
	\rule{0pt}{15pt}  $S$ \rule{0pt}{15pt} \\ $I$ \\ $R$ 
	\end{tabular} 
	&
	\begin{tabular}{r}
	\rule{0pt}{15pt}  45400 \\ 2100 \\ 2500
	\end{tabular}
	&
	\begin{tabular}{r}
	\rule{0pt}{15pt}  45737.1 \\ 1820.3 \\ 2442.6
	\end{tabular}
	&
	\begin{tabular}{r}
	\rule{0pt}{15pt}  337.1 \\ --279.7 \\ --57.4
	\end{tabular}
\end{tabular}
\end{center}

Do you see why there are differences?  We went forward in time using
the rates that were current at $t = 0$, but when we returned we used
the rates that were current at $t = 1$.  Because these rates were
different, we didn't get back where we started.  These differences do
not point to a flaw in the model; the problem lies with the way we are
trying to extract information from the model.  As we have been making
estimates,\index{estimate}
%
\mar{The differences measure how rough an estimate is}
%
we have assumed that the rates don't change over the course of a whole
day.  We already know that's not true, so the values that we have been
getting are not exact.  What this test adds to our knowledge is a way
to measure just {\em how\/} inexact those values are---as we do in the
table above.

\enlargethispage*{10pt}
In chapter 2 we will solve the problem of rough estimates by
recalculating all the quantities ten times a day, a hundred times a
day, or even more.  When we do the computations with shorter and
shorter time steps we will be able to see how the estimates improve.
We will even be able to see how to get values that are mathematically
exact!\label{`analyzeend'}

\medskip

\noindent
{\bf Delta notation}.\index{delta notation} This work has given us
some insights about the way our model predicts future values of $S$,
$I$, and $R$.  The basic idea is very simple: {\em determine how $S$,
  $I$, and $R$ change}.  Because these changes play such an important
role in what we do, it is worth having a simple way to refer to them.
Here is the notation that we will use:
\begin{center}
$\Delta x$ stands for a {\bf change}\index{change} in the quantity $x$ 
\end{center}
The symbol ``$\Delta$'' is the Greek capital letter {\em delta\/}; it
corresponds to the Roman letter ``D'' and stands for {\bf difference}.

Delta notation gives us a way to refer to changes of all sorts.  For
example, in the table on page~\pageref{`threeday'}, between day 1 and
day 3 the quantities $t$ and $S$ change by
\begin{align*}
	\Delta t &= 2 \text{ days},	\\
	\Delta S &=  -3010.9 \text{ persons}.
\end{align*}
We sometimes refer to a change as a {\bf step}.  
%
\mar{$\Delta$ stands for a change, a difference, or a step}
%
For instance, in this example we can say there is a ``$t$ step'' of 2
days, and an ``$S$ step'' of $-3010.9$ persons.  In the calculations
that produced the table on page~\pageref{`threeday'} we ``stepped
into'' the future, a day at a time.  Finally, delta notation gives us
a concise and vivid way to describe the relation between rates and
changes.  For example, if $S$ changes at the constant rate $S'$, then
under a $t$ step of $\Delta t$, the value of $S$ changes by
\[
\Delta S = S' \cdot \Delta t.
\]

\medskip
\noindent
{\bf Using the computer as a tool}.  Suppose we wanted to find out
what happens to $S$, $I$, and $R$ after a month, or even a year.  We
need only repeat---30 times, or 365 times---the three rounds of
calculations we used to go three days into the future.  The
computations would take time, though.  The same is true if we wanted
to do ten or one hundred rounds of calculations per day---which is the
approach we'll take in chapter 2 to get more accurate values.  To save
our time and effort we will soon begin to use a computer to do the
repetitive calculations.

\enlargethispage*{10pt}
A computer does calculations by following a set of instructions called
a {\bf program}.\index{program} Of course, if we had to give a million
instructions to make the computer carry out a million steps, there
would be no savings in labor.  The trick is to write a program with
just a few instructions that can be repeated over and over again to do
all the calculations we want.  The usual way to do this is to arrange
the instructions in a {\bf loop}\index{loop}\label{`loop'}.  To give
you an idea what a loop is, we'll look at the $S$-$I$-$R$
calculations.  They form a loop.  We can see the loop by making a flow
chart.

{\bf The flow chart}.\index{flow chart} We'll start by writing down
the three steps that take us from one day to the next:


\begin{description}
\item[Step I] Given the current values of $S$, $I$, and $R$, we get
  current $S'$, $I'$, and $R'$ by using the rate equations
\begin{alignat*}{2}
		S' &= - a S I, &   &     \\
		I' &= \hphantom{-}a S I & -\, & b I,  \\
		R' &=          &   & b I.
\end{alignat*}
\item[Step II] Given the current values of $S'$, $I'$, and $R'$, we
  find the changes $\Delta S$, $\Delta I$, and $\Delta R$ over the
  course of a day by using the equations
\begin{align*}
\Delta S \mbox{ persons} &= S' \ \dfrac{\mbox{persons}}{\mbox{day}} 
          \times 1 \ \mbox{day}, \\[4pt]
\Delta I \mbox{ persons} &= I' \ \dfrac{\mbox{persons}}{\mbox{day}} 
          \times 1 \ \mbox{day}, \\[4pt]
\Delta R \mbox{ persons} &= R' \ \dfrac{\mbox{persons}}{\mbox{day}} 
          \times 1 \ \mbox{day}.
\end{align*}
\item[Step III] Given the current values of $\Delta S$, $\Delta I$,
  and $\Delta R$, we find the new values of $S$, $I$, and $R$ a day
  later by using the equations
\begin{align*}
	\mbox{new $S$} &= \mbox{current $S$} + \Delta S, \\
	\mbox{new $I$} &= \mbox{current $I$} + \Delta I, \\
	\mbox{new $R$} &= \mbox{current $R$} + \Delta R.
\end{align*}
\end{description}

\noindent
Each step takes three numbers as {\bf input},\index{input} and
produces three new numbers as {\bf output}.\index{output} Note that
the output of each step is the input of the next, so the steps flow
together.  The diagram below, called a {\bf flow chart},\index{flow
  chart} 
%
\mar{\vspace{40pt} The flow chart\\ forms a loop} 
%
shows us how the steps are connected.

\begin{picture}(340,110)(5,0)\label{`flowchart'}
\put(0,50)
	{\begin{picture}(80,50)
	\put(0,5){\framebox(80,40)
		{\shortstack{values of\\[2pt]$S$, $I$, $R$}}}
	\end{picture}}
\put(130,50)
	{\begin{picture}(80,50)
	\put(0,5){\framebox(80,40)
		{\shortstack{values of\\[2pt]$S'$, $I'$, $R'$}}}
	\end{picture}}
\put(260,50)
	{\begin{picture}(80,50)
	\put(0,5){\framebox(80,40)
		{\shortstack{values of\\[2pt]
			$\Delta S$, $\Delta I$, $\Delta R$}}}
	\end{picture}}
\put(90,75){\vector(1,0){30}}
\put(220,75){\vector(1,0){30}}
\put(300,15){\line(0,1){30}}
\put(40,15){\vector(0,1){30}}
\put(40,15){\line(1,0){260}}
\put(105,85){\makebox[0pt]{I}}
\put(235,85){\makebox[0pt]{II}}
\put(170,25){\makebox[0pt]{III}}
\end{picture}

\noindent The calculations form a {\bf loop}\index{loop}, because the
output of step III is the input for step I.  If we go once around the
loop, the output of step III gives us the values of $S$, $I$, and $R$
{\em on the following day}.  The steps do indeed carry us into the
future.

Each step involves calculating three numbers.  If we count each
calculation as a single instruction, then it takes nine instructions
to carry the values of $S$, $I$, and $R$ one day into the future.  To
go a million days into the future, we need add only one more
instruction: ``Go around the loop a million times.''  In this way, a
computer program with only ten instructions can carry out a million
rounds of calculations!

Later in this chapter (section~3) you will find a real computer
program that lists these instructions (for three days instead of a
million, though).
%
\mar{A computer program\\ will carry out\\ the three steps}
%
Study the program to see which instructions accomplish which steps.
In particular, see how it makes a loop.  Then run the program to check
that the computer reproduces the values you already computed by hand.
Once you see how the program works, you can modify it to get further
information---for example, you can find out what happens to $S$, $I$,
and $R$ thirty days into the future.  You will even be able to plot
the graphs of $S$, $I$, and $R$.

\aside{Rate equations\index{rate equation} have always been at the
  heart of calculus, and they have been analyzed using mechanical and
  electronic computers for as long as those tools have been available.
  Now that small powerful computers have begun to appear in the
  classroom, it is possible for beginning calculus students to explore
  interesting and complex problems that are modelled by rate
  equations.  Computers are changing how mathematics is done and how
  it is learned.}

\medskip

\noindent {\bf Analysis without a computer}.  A computer is a powerful
tool for exploring the $S$-$I$-$R$ model, but there are many things we
can learn about the model without using a computer.  Here is an
example.

According to the model, the rate at which the infected population
grows is given by the equation
\[
I' = .00001 \, S I - I/14 \quad \mbox{persons/day.}
\]
In our example, $I' = 803.4$ at the outset.  This is a positive
number, so $I$ increases initially.  In fact, $I$ will continue to
increase as long as $I'$ is positive.  If $I'$ ever becomes negative,
then $I$ decreases.  So let's ask the question: when is $I'$ positive,
when is it negative, and when is it zero?  By factoring out $I$ in the
last equation we obtain
\[
I' = I \left( .00001 \, S - \dfrac{1}{14} \right) \ \mbox{persons/day.}
\]
Consequently $I' = 0$ if either
\[
I = 0 \qquad \mbox{or} \qquad .00001 \, S - \dfrac{1}{14} = 0.
\]
The first possibility $I = 0$ has a simple interpretation: there is no
infection within the population.  The second possibility is more
interesting; it says that $I'$ will be zero when
\[
.00001 \, S - \dfrac{1}{14} = 0 \qquad \mbox{or} \qquad
		S = \dfrac{100000}{14} \approx 7142.9.
\]

If $S$ is {\em greater\/} than $100000/14$ and $I$ is positive, then
you can check that the formula
\[
I' = I \left( .00001 \, S - \dfrac{1}{14} \right) 
		\ \mbox{persons/day}
\]
tells us $I'$ is positive---so $I$ is increasing.  If, on the other
hand, $S$ is {\em less than\/} $100000/14$, then $I'$ is negative and
$I$ is decreasing.  So $S = 100000/14$ represents a {\bf
  threshold}.\index{threshold}\label{`threshold'} If $S$ falls below
the threshold, $I$ decreases.  If $S$ exceeds the threshold, $I$
increases.  Finally, $I$ reaches its peak when $S$ equals the
threshold.

The presence of a threshold value for $S$ is purely a  mathematical result.  However, it has an interesting interpretation for the epidemic.  
%
\mar{The threshold determines whether there will be an epidemic}
%
As long as there are at least $7143$ susceptibles, the infection will
spread, in the sense that there will be more people falling ill than
recovering each day.  As new people fall ill, the number of
susceptibles declines.  Finally, when there are fewer than $7143$
susceptibles, the pattern reverses: each day more people will recover
than will fall ill.

If there were fewer than $7143$ susceptibles in the population {\em at
  the outset}, then the number of infected would only decline with
each passing day.  The infection would simply never catch hold.  The
clear implication is that the noticeable surge in the number of cases
that we associate with an ``epidemic'' disease is due to the presence
of a large susceptible population.  If the susceptible population lies
below a certain threshold value, a surge just isn't possible.  This is
a valuable insight---and we got it with little effort.  We didn't need
to make lengthy calculations or call on the resources of a computer; a
bit of algebra was enough.

\newpage

\subsection{Exercises}

\subsubsection{Reading a Graph}

\index{graph} The graphs on pages \pageref{Sgraph}
and~\pageref{Rgraph} have no scales marked along their axes, so they
provide mainly {\em qualitative\/} information.  The graphs below do
have scales, so you can now answer {\em quantitative\/} questions
about them.  For example, on day 20 there are about 18,000 susceptible
people.  Read the graphs to answer the following questions.  (Note: $S
+ I + R$ is {\em not\/} constant in this example, so these graphs
cannot be solutions to our model..)

\setlength{\unitlength}{1pt}
%\begin{picture}(340,140)(15,0)
\begin{picture}(340,135)(15,0)
	\put(310,25){time}
	\put(56,114){$S$}
	\put(90,20){\line(0,1){90}}
	\put(130,20){\line(0,1){90}}
	\put(170,20){\line(0,1){90}}
	\put(210,20){\line(0,1){90}}
	\put(250,20){\line(0,1){90}}
	\put(90,8){\makebox[0pt]{10}}
	\put(130,8){\makebox[0pt]{20}}
	\put(170,8){\makebox[0pt]{30}}
	\put(210,8){\makebox[0pt]{40}}
	\put(250,8){\makebox[0pt]{50}}
	\put(290,8){days}
	\put(50,57){\line(1,0){230}}
	\put(50,94){\line(1,0){230}}
	\put(46,53){\makebox[0pt][r]{$20000$}}
	\put(46,90){\makebox[0pt][r]{$40000$}}
	\put(46,114){\makebox[0pt][r]{people}}
	\begin{thicklines}
	\put(50,20){\vector(1,0){255}}
	\put(50,20){\vector(0,1){100}}
	\bezier{75}(50,100)(70,94)(100,71)
	\bezier{50}(100,71)(160,25)(300,24) %200
	\end{thicklines}
\end{picture}

%\begin{picture}(340,140)(15,0)
\begin{picture}(340,130)(15,0)
	\put(310,25){time}
	\put(56,114){$I$}
	\put(90,20){\line(0,1){90}}
	\put(130,20){\line(0,1){90}}
	\put(170,20){\line(0,1){90}}
	\put(210,20){\line(0,1){90}}
	\put(250,20){\line(0,1){90}}
	\put(90,8){\makebox[0pt]{10}}
	\put(130,8){\makebox[0pt]{20}}
	\put(170,8){\makebox[0pt]{30}}
	\put(210,8){\makebox[0pt]{40}}
	\put(250,8){\makebox[0pt]{50}}
	\put(290,8){days}
	\put(50,47){\line(1,0){230}}
	\put(50,74){\line(1,0){230}}
	\put(50,101){\line(1,0){230}}
	\put(46,44){\makebox[0pt][r]{$5000$}}
	\put(46,71){\makebox[0pt][r]{$10000$}}
	\put(46,98){\makebox[0pt][r]{$15000$}}
	\put(46,116){\makebox[0pt][r]{people}}
	\begin{thicklines}
	\put(50,20){\vector(1,0){255}}
	\put(50,20){\vector(0,1){100}}
	\bezier{75}(50,25)(65,35)(75,60)
	\bezier{50}(75,60)(99,120)(124,72) %120
	\bezier{50}(124,72)(149,24)(300,23) %200
	\end{thicklines}
\end{picture}


%\begin{picture}(340,140)(15,0)
\begin{picture}(340,130)(15,0)
	\put(310,25){time}
	\put(56,120){$R$}
	\put(90,20){\line(0,1){90}}
	\put(130,20){\line(0,1){90}}
	\put(170,20){\line(0,1){90}}
	\put(210,20){\line(0,1){90}}
	\put(250,20){\line(0,1){90}}
	\put(90,8){\makebox[0pt]{20}}
	\put(130,8){\makebox[0pt]{40}}
	\put(170,8){\makebox[0pt]{60}}
	\put(210,8){\makebox[0pt]{80}}
	\put(250,8){\makebox[0pt]{100}}
	\put(290,8){days}
	\put(50,57){\line(1,0){230}}
	\put(50,94){\line(1,0){230}}
	\put(46,53){\makebox[0pt][r]{$20000$}}
	\put(46,90){\makebox[0pt][r]{$40000$}}
	\put(46,120){\makebox[0pt][r]{people}}
	\begin{thicklines}
	\put(50,20){\vector(1,0){255}}
	\put(50,20){\vector(0,1){100}}
	\bezier{90}(50,30)(90,40)(120,70)
	\bezier{50}(120,70)(160,110)(300,110) %200
	\end{thicklines}
\end{picture}

\vspace{-10pt} 

\enlargethispage*{10pt}
\num When does the infection hit its peak?\index{maximum} How many
people are infected at that time?

\num Initially, how many people are susceptible?  How many days does
it take for the susceptible population to be cut in half?

\num How many days does it take for the recovered population to reach
25,000?  How many people {\em eventually\/} recover?  Where did you
look to get this information?

\num On what day is the size of the infected population increasing
most rapidly?  When is it decreasing most rapidly?  How do you know?

\num How many people caught the illness at some time during the first
20 days?  (Note that this is not the same as the number of people who
are infected on day 20.)  Explain where you found this information.

\num Copy the graph of $R$ as accurately as you can, and then
superimpose a sketch of $S$ on it.  Notice the time scales on the {\em
  original\/} graphs of $S$ and $R$ are different.  Describe what
happened to the graph of $S$ when you superimposed it on the graph of
$R$.  Did it get compressed or stretched?  Was this change in the
horizontal direction or the vertical?

\subsubsection{A Simple Model}

These questions concern the rate equation $S' = -470$ persons per day
that we used to model a susceptible population on
pages \pageref{`modelstart'}--\pageref{`modelend'}.

\num Suppose the initial susceptible population was 20,000 on
Wednesday.  Use the model to answer the following questions.

\sub  How many susceptibles will be left ten days later?

\sub How many days will it take for the susceptible population to
vanish entirely?

\sub  How many susceptibles were there on the previous Sunday?  

\sub How many days before Wednesday were there 30,000 susceptibles?



\subsubsection{Mark Twain's Mississippi}
\index{Mississippi}

The Lower Mississippi River meanders over its flat valley, forming
broad loops called ox-bows.\index{ox-bow} In a flood, the river can
jump its banks and cut off one of these loops, getting shorter in the
process.  In his book {\em Life on the Mississippi} (1884), Mark
Twain\index{Twain, M.} suggests, with tongue in cheek, that some day
the river might even vanish!  Here is a passage that shows us some of
the pitfalls in using rates to predict the future and the past.
\begin{quote}\label{`twain'}
{\small In the space of one hundred and seventy six years the Lower
  Mississippi has shortened itself two hundred and forty-two miles.
  That is an average of a trifle over a mile and a third per year.
  Therefore, any calm person, who is not blind or idiotic, can see
  that in the Old O\"{o}litic Silurian Period, just a million years
  ago next November, the Lower Mississippi was upwards of one million
  three hundred thousand miles long, and stuck out over the Gulf of
  Mexico like a fishing-pole.  And by the same token any person can
  see that seven hundred and forty-two years from now the Lower
  Mississippi will be only a mile and three-quarters long, and Cairo
  [Illinois] and New Orleans will have joined their streets together
  and be plodding comfortably along under a single mayor and a mutual
  board of aldermen.  There is something fascinating about science.
  One gets such wholesome returns of conjecture out of such a trifling
  investment of fact.}
\end{quote}

\noindent  
Let $L$ be the length of the Lower Mississippi River.  Then $L$ is a
variable quantity we shall analyze.

\num According to Twain's data, what is the exact {\bf
  rate}\index{rate} at which $L$ is changing, in miles per year?  What
approximation does he use for this rate?  Is this a reasonable
approximation? Is this rate {\em positive\/} or {\em negative\/}?
Explain.  In what follows, use Twain's approximation.

\num Twain wrote his book in 1884.  Suppose the Mississippi that Twain
wrote about had been 1100 miles long; how long would it have become in
1990?

\num Twain does not tell us how long the Lower Mississippi was in 1884
when he wrote the book, but he does say that 742 years later it will
be only $1\tfrac{3}{4}$ miles long.  How long must the river have been
when he wrote the book?

\num Suppose $t$ is the number of years since 1884.  Write a formula
that describes how much $L$ has changed in $t$ years.  Your formula
should complete the equation
\[
\mbox{the change in $L$ in $t$ years} = \ldots \ . \rule{90pt}{0pt}
\]

\num From your answer to question 10, you know how long the river was
in 1884.  From question 11, you know how much the length has changed
$t$ years after 1884.  Now write a formula that describes how long the
river is $t$ years later.

\num Use your formula to find what $L$ was a million years ago.  Does
your answer confirm Twain's assertion that the river was ``upwards of
1,300,000 miles long'' then?

\num Was the river ever 1,300,000 miles long; will it ever be
$1\tfrac{3}{4}$ miles long?  (This is called a {\bf reality check}.)
What, if anything, is wrong with the ``trifling investment of fact''
which led to such ``wholesale returns of conjecture'' that Twain has
given us?


\subsubsection{The Measles Epidemic}

We consider once again the specific rate equations 
\begin{align*}
S' &= -.00001 \, SI, \\
I' &= .00001 \, SI - I/14, \\
R' &= I/14, 
\end{align*}
discussed in the text on pages
\pageref{`analyzestart'}--\pageref{`analyzeend'}.  We saw that at time
$t = 1$,
\[
S = 44446.6, \qquad I = 2903.4, \qquad R = 2650.0.
\]

\num Calculate the current rates of change $S'$, $I'$, and $R'$ when
$t = 1$, and then use these values to determine $S$, $I$, and $R$ one
day later\label{`measles'}.

\num In the previous question you found $S$, $I$, and $R$ when $t =
2$.  Using these values, calculate the rates $S'$, $I'$, and $R'$ and
then determine the new values of $S$, $I$, and $R$ when $t = 3$.  See
the table on page~\pageref{`threeday'}.

\num {\bf Double the time step}.  Go back to the starting time $t = 0$
and to the initial values
\[
S = 45400, \qquad I = 2100, \qquad R = 2500.
\]
Recalculate the values of $S$, $I$, and $R$ at time $t = 2$ by using a
time step of $\Delta t = 2$.  You should perform only a single round
of calculations, and use the rates $S'$, $I'$, and $R'$ that are
current at time $t = 0$.


\num {\bf There and back again}.\index{there and back
  again}\label{`therebackprob'} In the text we went one day into the
future and then back again to the present.  Here you'll go forward two
days from $t = 0$ and then back again.  There are two ways to do this:
with a time step of $\Delta t = \pm 2$ (as in the previous question),
and with a pair of time steps of $\Delta t = \pm 1$ .

\sub ($\Delta t = \pm 2$).  Using the values of $S$, $I$, and $R$ at
time $t = 2$ that you just got in the previous question, calculate the
rates $S'$, $I'$, and $R'$.  Then using a time step of $\Delta t =
-2$, estimate new values of $S$, $I$, and $R$ at time $t = 0$.  How
much do these new values differ from the original values 45,400, 2100,
2500?

\sub  ($\Delta t = \pm 1$).  Now make a new start, using the values 
\begin{xalignat*}{3}
S  &= 43156.1, &  I  &= 3986.5, &  R  &= 2857.4, \\ 
S' &= -1720.4, &  I' &= 1435.7, &  R' &= 284.7.
\end{xalignat*}
that occur when $t = 2$ if we make estimates with a time step $\Delta
t = 1$.  (These values come from the table on
page~\pageref{`threeday'}) Using two rounds of calculations with a
time step of $\Delta t = -1$, estimate another set of new values for
$S$, $I$, and $R$ at time $t = 0$.  How much do these new values
differ from the original values 45,400, 2100, 2500?

\sub Which process leads to a {\em smaller\/} set of differences: a
single round of calculations with $\Delta t = \pm 2$, or two rounds of
calculations with $\Delta t = \pm 1$?  Consequently, which process
produces better estimates---in the sense in which we used to measure
estimates on page~\pageref{`differences'}?


\num {\bf Quarantine}.\index{quarantine} One of the ways to treat an
epidemic is to keep the infected away from the susceptible; this is
called quarantine.  The intention is to reduce the chance that the
illness will be transmitted to a susceptible person.  Thus, quarantine
alters the {\em transmission coefficient}.

\sub Suppose a quarantine is put into effect that cuts in half the
chance that a susceptible will fall ill.  What is the new transmission
coefficient?

\sub On page~\pageref{`threshold'} it was determined that whenever
there were fewer than 7143 susceptibles, the number of infected would
decline instead of grow.  We called 7143 a {\em
  threshold\/}\index{threshold} level for $S$.  Changing the
transmission coefficient, as in part (a), changes the threshold level
for $S$.  What is the new threshold?

\sub Suppose we start with $S$ = 45,400.  Does quarantine eliminate
the epidemic, in the sense that the number of infected immediately
goes down from 2100, without ever showing an increase in the number of
cases?

\sub Since the new transmission coefficient is not small enough to
guarantee that $I$ never goes up, can you find a smaller value that
{\em does\/} guarantee $I$ never goes up?  Continue to assume we start
with $S = 45400$.

\sub Suppose the initial susceptible population is 45,400.  What is
the {\em largest\/} value that the transmission coefficient can have
and still guarantee that $I$ never goes up?  What level of quarantine
does this represent?  That is, do you have to reduce the chance that a
susceptible will fall ill to one-third of what it was with no
quarantine at all, to one-fourth, or what?


\subsubsection{Other Diseases}

\vspace{-10pt}

\num Suppose the spread of an illness similar to measles is modelled
by the following rate equations:
\begin{align*}
S' &= -.00002 \, SI, \\
I' &= .00002 \, SI - .08 \, I, \\
R' &= .08 \, I.
\end{align*}
Note: the initial values $S = 45400$, etc.\ that we used in the text
do not apply here.

\sub Roughly how long does someone who catches this illness remain
infected?  Explain your reasoning.

\sub How large does the susceptible population have to be in order for
the illness to take hold---that is, for the number of cases to
increase?  Explain your reasoning.

\sub Suppose 100 people in the population are currently ill.
According to the model, how many (of the 100 infected) will recover
during the next 24 hours?

\sub Suppose 30 {\em new\/} cases appear during the same 24 hours.
What does that tell us about $S'$?

\sub Using the information in parts (c) and (d), can you determine how
large the current susceptible population is?

\numa Construct the appropriate $S$-$I$-$R$ model for a measles-like
illness that lasts for 4 days.  It is also known that a typical
susceptible person meets only about 0.3\% of infected population each
day, and the infection is transmitted in only one contact out of six.

\sub How small does the susceptible population have to be for this
illness to fade away without becoming an epidemic?

\num Consider the general $S$-$I$-$R$ model for a measles-like
illness:
\begin{alignat*}{2}
S' &=  - a S I,         &     &     \\
I' &= \hphantom{-}a S I & -\, & b I,  \\
R' &=                   &     & b I.
\end{alignat*}

\sub The threshold level for $S$---below which the number of infected
will only decline---can be expressed in terms of the transmission
coefficient $a$ and the recovery coefficient $b$.  What is that
expression?\label{`setthresh'}

\sub Consider two illnesses with the same transmission coefficient
$a$; assume they differ only in the length of time someone stays ill.
Which one has the lower threshold level for $S$?  Explain your
reasoning.


\subsubsection{What Goes Around Comes Around}

Some relatively mild illnesses, like the common cold, return to infect
you again and again.  For a while, right after you recover from a
cold, you are immune.  But that doesn't last; after some weeks or
months, depending on the illness, you become susceptible again.  This
means there is now a flow from the recovered population to the
susceptible.  These exercises ask you to modify the basic $S$-$I$-$R$
model to describe an illness where immunity is temporary.

\num Draw a compartment diagram for such an illness.  Besides having
all the ingredients of the diagram on page~\pageref{`diagram'}, it
should depict a flow from $R$ to~$S$.  Call this {\bf immunity
  loss},\index{immunity loss} and use $c$ to denote the coefficient of
immunity loss.

\num Suppose immunity is lost after about six weeks.  Show that you
can set $c = 1/42$ per day, and explain your reasoning carefully.  A
suggestion: adapt the discussion of recovery in the text.

\num Suppose this illness lasts 5 days and it has a transmission
coefficient of .00004 in the population we are considering.  Suppose
furthermore that the total population is fixed in size (as was the
case in the text).  Write down rate equations for $S$, $I$, and $R$.

\num We saw in the text that the model for an illness that confers
permanent immunity has a threshold value for $S$ in the sense that
when $S$ is above the threshold, $I$ increases, but when it is below,
$I$ decreases.  Does {\em this\/} model have the same feature?  If so,
what is the threshold value?

\num For a mild illness that confers permanent immunity, the size of
the recovered population can only grow.  This question explores what
happens when immunity is only temporary.

\sub  Will $R$ increase or decrease if 
\[
S = 45400, \qquad I = 2100, \qquad R = 2500?
\]

\sub Suppose we shift 20000 susceptibles to the recovered population
(so that $S = 25400$ and $R = 22500$), leaving $I$ unchanged.  Now,
will $R$ increase or will it decrease?

\sub Using a total population of 50,000, give two other sets of values
for $S$, $I$, and $R$ that lead to a decreasing $R$.

\sub In fact, the relative sizes of $I$ and $R$ determine whether $R$
will increase or decrease.  Show that
\begin{align*}
\text{if } & I > \tfrac{5}{42} R, \quad 
		\text{then $R$ will increase;} \\
\text{if } & I < \tfrac{5}{42} R,  \quad
		\text{then $R$ will decrease.} 
\end{align*}
Explain your argument clearly.  A suggestion: consider the rate
equation for~$R'$.

\num {\bf The steady state}.  Any illness that confers only temporary
immunity can appear to linger in a population forever.  You may not
always have a cold, but someone does, and eventually you catch another
one.  (``What goes around comes around.'')  Individuals gradually move
from one compartment to the next.  When they return to where they
started, they begin another cycle.

Each compartment (in the diagram you drew in exercise~23) has an {\em
  inflow\/} and an {\em outflow}.  It is conceivable that the two
exactly balance, so that the {\em size\/} of the compartment doesn't
change (even though its individual occupants do).  When this happens
for all three compartments simultaneously, the illness is said to be
in a {\bf steady state}.  In this question you explore the steady
state of the model we are considering.  Recall that the total
population is 50,000.

\sub What must be true if the inflow and outflow to the $I$
compartment are to balance?

\sub What must be true if the inflow and outflow to the $R$
compartment are to balance?

\sub If neither $I$ nor $R$ is changing, then the model must be at the
steady state.  Why?

\sub  What is the value of $S$ at the steady state?

\sub What is the value of $R$ at the steady state?  A suggestion: you
know $R + I = 50000 - \mbox{(the steady state value of $S$)}$.  You
also have a connection between $I$ and $R$ at the steady state.


%%%%%%%%%%%%%%%%%%%%%%%%%%%%%%%%%%%%%%%%%%%%%%%%%%%%%%%%%%%%%%%%%%%%%%
%                                                                    %
%                             Section 2                              %
%                                                                    %  
%%%%%%%%%%%%%%%%%%%%%%%%%%%%%%%%%%%%%%%%%%%%%%%%%%%%%%%%%%%%%%%%%%%%%%


\section{The Mathematical Ideas}

A number of important mathematical ideas have already emerged in our
study of an epidemic.  In this section we pause to consider them,
because they have a ``universal'' character.  Our aim is to get a
fuller understanding of what we have done so we can use the ideas in
other contexts.

\aside{We often draw out of a few particular experiences a lesson that
  can be put to good use in new settings.  This process is the essence
  of mathematics, and it has been given a name---abstraction---which
  means literally ``drawing from.''\index{abstraction} Of course
  abstraction is not unique to mathematics; it is a basic part of the
  human psyche.}

\vspace{-13pt}

\subsection{Functions}

\index{function}

A {\bf function}\index{function} describes how one quantity depends on
another.  In our study of a measles epidemic,
%
\mar{Functions and\\ their notation}
%
the relation between the number of susceptibles $S$ and the time $t$
is a function.  We write $S(t)$ to denote that $S$ is a function of
$t$.  We can also write $I(t)$ and $R(t)$ because $I$ and $R$ are
functions of $t$, too.  We can even write $S'(t)$ to indicate that the
rate $S'$ at which $S$ changes over time is a function of $t$.  In
speaking, we express $S(t)$ as ``$S$ of $t$'' and $S'(t)$ as ``$S$
prime of $t$.''

You can find functions everywhere.  The amount of postage you pay for
a letter is a function of the weight of the letter.  The time of
sunrise is a function of what day of the year it is.  The crop yield
from an acre of land is a function of the amount of fertilizer used.
The position of a car's gasoline gauge (measured in centimeters from
the left edge of the gauge) is a function of the amount of gasoline in
the fuel tank.  On a polygraph (``lie detector'') there is a pen that
records breathing; its position is a function of the amount of
expansion of the lungs.  The volume of a cubical box is a function of
the length of a side.  The last is a rather special kind of function
because it can be described by an algebraic formula: if $V$ is the
volume of the box and $s$ is the length of a side, then $V(s) = s^3$.

Most functions are {\em not\/} described by algebraic formulas,
however.\index{formula} For instance, the postage function is given by
a set of verbal instructions and the time of sunrise is given by a
table in an almanac.  The relation between a gas gauge and the amount
of fuel in the tank is determined simply by making measurements.
There is no algebraic formula that tells us how the number of
susceptibles, $S$, depends upon $t$, either.  Instead, we find $S(t)$
by carrying out the steps in the flow chart on
page~\pageref{`flowchart'} until we reach $t$ days into the future.

In the function $S(t)$ the variable $t$ is called the {\bf
  input}\index{input}
%
\mar{A function has input and output}
%
and the variable $S$ is called the {\bf output}.\index{output} In the
sunrise function, the day of the year is the input and the time of
sunrise is the output.  In the function $S(t)$ we think of $S$ as {\em
  depending on\/} $t$, so $t$ is also called the independent variable
and $S$ the dependent variable.  The set of values that the input
takes is called the {\bf domain}\index{domain} of the function.  The
set of values that the output takes is called the {\bf
  range}.\index{range}

The idea of a function is one of the central notions of mathematics.
It is worth highlighting:
\begin{center}
\begin{picture}(330,60)
	\begin{thicklines}
	\put(0,0){\framebox(330,60)
		{\shortstack{\bf A function is a rule that 
			specifies how \\
		\bf the value of one variable, the input, 
			determines \\
		\bf the value of a second variable, the output.}}}
	\end{thicklines}
\end{picture}
\end{center}
Notice that we say {\em rule\/} here, and not {\em formula}.  This is
deliberate.  We want the study of functions to be as broad as
possible, to include all the ways one quantity is likely to be related
to another in a scientific question.

\medskip
\noindent
{\bf Some technical details}.  It is important not to confuse an
expression like $S(t)$ with a product; $S(t)$ does {\em not\/} mean $S
\times t$.  On the contrary, the expression $S(1.4)$, for example,
stands for the output of the function $S$ when 1.4 is the input.  In
the epidemic model we then interpret it as the number of susceptibles
that remain 1.4 days after today.

We have followed the standard practice in science by letting the
single letter $S$ designate both the {\em function\/}---that is, the
{\em rule\/}---and the {\em output\/} of that function.  Sometimes,
though, we will want to make the distinction.  In that case we will
use two different symbols.  For instance, we might write $S = f(t)$.
Then we are still using $S$ to denote the output, but the new symbol
$f$ stands for the function rule.

The symbols we use to denote the input and the output of a function
are just names; if we change them, we don't change the function.  For
example, here are three ways to describe the same function $g$:
\begin{align*}
g &: \mbox{multiply the input by 5, then subtract 3} \\
g(x) &= 5x - 3 \\
g(u) &= 5u - 3 \, .
\end{align*}
It is important to realize that the {\em formulas} we just wrote in
the last two lines are merely shorthand for the instructions stated in
the first line.  If you keep this in mind, then absurd-looking
combinations like $g(g(2))$ can be decoded easily by remembering $g$
of {\em anything} is just 5 times that thing, minus 3.  We could thus
evaluate $g(g(2))$ from the inside out (which is usually easier) as
$$
g(g(2))=g(5\cdot 2-3)=g(7)=5\cdot 7 -3=32,
$$
or we could evaluate it from the outside in as  
$$
g(g(2))=5g(2)-3=5(5\cdot 2 -3)-3=5\cdot 7-3=32,
$$ 
as before. 

Suppose $f$ is some other rule, say $f(t)=t^2+4t-1$.  Remember that
this is just shorthand for ``Take the input (whatever it is), square
it, add four times the input, and subtract 1.''  We could then
evaluate
\[
f(g(3))=f(5\cdot 3 -3)=f(12)= 12^2+4\cdot 12 - 1 = 144+48-1=191,
\]
while 
\[
g(f(3))=g(20)=97.
\] 

This process of {\bf chaining} functions together by using the output
of one function as the input for another turns out to be very
\mar{Part of math is simply learning the language} important later in
this course, and will be taken up again in chapter 3.  For now,
though, you should treat it simply as part of the formal language of
mathematics, requiring a knowledge of the rules but no cleverness.  It
is analogous to learning how to conjugate verbs in French class---it's
not very exciting for its own sake, but it allows us to read the
interesting stuff later on!

A particularly important class of functions is composed of the {\bf
  constant functions} which give the same output for every input.  If
$h$ is the constant function that always gives back 17, then in
formula form we would express this as $h(x) = 17$.  Constant functions
are so simple you might feel you are missing the point, but that's all
there is to it!

\subsection{Graphs}
\index{graph}

A graph describes a function in a visual form.  Sometimes---as with a
seismograph or a lie detector,
%
\mar{A graph is a function rule given visually}
%
for instance---this is the {\em only\/} description we have of a
particular function.  The usual arrangement is to put the input
variable on the horizontal axis and the output on the vertical---but
it is a good idea when you are looking at a particular graph to take a
moment to check; sometimes, the opposite convention is used!  This is
often the case in geology and economics, for instance.


\setlength{\unitlength}{1pt}
\begin{picture}(340,140)(15,0)
	\put(50,20){\vector(1,0){255}}
	\put(50,20){\vector(0,1){100}}
	\put(310,25){$t$}
	\put(56,114){$S$}
	\put(90,20){\line(0,1){3}}
	\put(130,20){\line(0,1){3}}
	\put(170,20){\line(0,1){3}}
	\put(90,8){\makebox[0pt]{10}}
	\put(130,8){\makebox[0pt]{20}}
	\put(170,8){\makebox[0pt]{30}}
	\put(290,8){days}
	\put(50,57){\line(1,0){3}}
	\put(50,94){\line(1,0){3}}
	\put(46,53){\makebox[0pt][r]{$20000$}}
	\put(46,90){\makebox[0pt][r]{$40000$}}
	\put(46,114){\makebox[0pt][r]{people}}
	\put(100,20){\circle*{3}}
	\put(100,71){\circle*{3}}
	\put(50,71){\circle*{3}}
	\put(100,6){\makebox[0pt][l]{$t_0$}}
	\put(46,70){\makebox[0pt][r]{$S_0$}}
	\put(104,73){\makebox[0pt][l]{$(t_0, S_0)$}}
	\multiput(100,71)(0,-5){10}{\circle*{1}}
	\multiput(100,71)(-5,0){10}{\circle*{1}}
	\begin{thicklines}
	\bezier{75}(50,100)(70,94)(100,71)
	\bezier{200}(100,71)(160,25)(300,24)
	\end{thicklines}
\end{picture}

Sketched above is the graph of a function $S(t)$ that tells how many
susceptibles there are after $t$ days.  Given any $t_0$, we ``read''
the graph to find $S(t_0)$, as follows: from the point $t_0$ on the
$t$-axis, go vertically until you reach the graph; then go
horizontally until you reach the $S$-axis.  The value $S_0$ at that
point is the output $S(t_0)$.  Here $t_0$ is about 13 and $S_0$ is
about 27,000; thus, the graph says that $S(13) \approx 27000$, or
about 27,000 susceptibles are left after 13 days.


\subsection{Linear Functions}
\index{linear function}

{\bf Changes in input and output}.  
%
\mar{If $y$ depends on $x$, then $\Delta y$ depends on $\Delta x$}
%
Suppose $y$ is a function of $x$.  Then there is some rule that
answers the question: What is the value of $y$ for any given $x$?
Often, however, we start by knowing the value of $y$ for a particular
$x$, and the question we really want to ask is: How does $y$ respond
to {\em changes\/} in $x$?  We are still dealing with the same
function---just looking at it from a different point of view.  This
point of view is important; we use it to analyze functions (like
$S(t)$, $I(t)$, and $R(t)$) that are defined by rate equations.

The way $\Delta y$ depends on $\Delta x$ can be simple or it can be
complex, depending on the function involved.  The simplest possibility
is that $\Delta y$ and $\Delta x$ are {\bf
  proportional}:\index{proportional}
\[
\Delta y = m\cdot\Delta x, \quad 
		\mbox{for some constant $m$.}
\]
Thus, if $\Delta x$ is doubled, so is $\Delta y$; if $\Delta x$ is
tripled, so is $\Delta y$.
%
\mar{\vspace{-48pt} The defining property of a linear function}
%
A function whose input and output are related in this simple way is
called a {\bf linear function}, because the graph is a straight line.
Let's take a moment to see why this is so.

\medskip
\noindent
\begin{minipage}[b]{3.5in}
{\bf The graph of a linear function}.\index{linear function!graph}
 The graph consists of certain points $(x, y)$ in the $x, y$-plane.
Our job is to see how those points are arranged.  Fix one of them, and
call it $(x_0, y_0)$.  Let $(x, y)$ be any other point on the graph.
Draw the line that connects this point to $(x_0, y_0)$, as we have
done in the figure at the right.  Now set
	\vspace{-3pt}
\[
\Delta x = x - x_0, \qquad \Delta y = y - y_0.
\]
\end{minipage}
\begin{minipage}[b]{1.3in}
\begin{picture}(170,100)(0,20)
	\put(20,25){\vector(1,0){160}}
	\put(35,20){\vector(0,1){115}}
	\put(170,14){\small $x$}
	\put(40,125){\small $y$}
	\begin{thicklines}
		\put(65,60){\line(3,2){66}}
	\end{thicklines}
	\put(77,68){\circle*{4}}
	\put(119,96){\circle*{4}}
	\put(79,74){\makebox[0pt][r]
		{\small $(x_0, y_0)$}}
	\put(121,102){\makebox[0pt][r]
		{\small $(x, y)$}}
	\put(77,68){\line(1,0){42}}
	\put(119,68){\line(0,1){28}}
	\put(100,58){\makebox[0pt]
		{\small $\Delta x$}}
	\put(122,78){\makebox[0pt][l]
		{\small $\Delta y=m\cdot\Delta x$}}
\end{picture}
\end{minipage}

\medskip
\noindent
By definition of a linear function, $\Delta y=m\cdot\Delta x$, as the
figure shows, so the slope of this line is $\Delta y/\Delta x = m$.
Recall that $m$ is a constant; thus, if we pick a new point $(x, y)$,
the slope of the connecting line won't change.

Since $(x, y)$ is an arbitrary point on the graph, what we have shown
is that {\bf every point on the graph lies on a line of slope
  \boldmath $m$ through the point \boldmath $(x_0, y_0)$}.  But there
is only one such line---and all the points lie on it!  That line must
be the graph of the linear function.


\begin{center}
\begin{picture}(330,45)
	\begin{thicklines}
	\put(0,0){\framebox(330,45)
		{\shortstack{\bf A linear function is one that 
			satisfies \boldmath 
			$\Delta y = m\cdot\Delta x$;\\
		 \bf its graph is a straight line whose slope is
			\boldmath $m$.}}}
	\end{thicklines}
\end{picture}
\end{center}


\noindent
{\bf Rates, slopes, and
  multipliers}.\index{rate}\index{slope}\index{multiplier}\label{multiplier-begin}
The interpretation of $m$ as a slope is just one possibility; there
are two other interpretations that are equally important.  To
illustrate them we'll use Mark Twain's vivid description of the
shortening of the Lower Mississippi River (see
page~\pageref{`twain'}).  This will also give us the chance to see how
a linear function emerges in context.

\enlargethispage*{10pt}
Twain says ``the Lower Mississippi has shortened itself \ldots\ an
average of a trifle over a mile and a third per year.''  Suppose we
let $L$ denote the length of the river, in miles, and $t$ the time, in
years.  Then $L$ depends on $t$, and Twain's statement implies that
$L$ is a {\em linear\/} function of $t$---in the sense in which we
have just defined a linear function.  Here is why.  According to our
definition, there must be some number $m$ which makes $\Delta L =
m\cdot\Delta t$.  But notice that Twain's statement has exactly this
form if we translate it into mathematical language.  Convince yourself
that it 
%
\mar{\vspace{12pt} Stop and do\\ the translation}
%
says 
\[
\Delta L \ \mbox{miles} = -1 \tfrac{1}{3} \ 
			\dfrac{\text{miles}}{\text{year}}
		\times \Delta t \ \mbox{years}.
\]
Thus we should take $m$ to be $-1\tfrac{1}{3}$ miles per year.  


The role of $m$ here is to convert one quantity ($\Delta t$ years)
into another ($\Delta L$ miles) by multiplication.  All linear
functions work this way.  In the defining equation $\Delta y =
m\cdot\Delta x$, multiplication by $m$ converts $\Delta x$ into
$\Delta y$.  Any change in $x$ produces a change in $y$ that is $m$
times as large.  For this reason we give $m$ its second interpretation
as a {\bf multiplier}.

\aside{It is easier to understand why the usual symbol for {\em
    slope\/} is $m$---instead of $s$---when you see that a slope can
  be interpreted as a multiplier.}

It is important to note that, in our example, $m$ is not simply
$-1\tfrac{1}{3}$; it is $-1\tfrac{1}{3}$ {\em miles per year}.  In
other words, $m$ is the {\bf rate} at which the river is getting
shorter.  All linear functions work this way, too.  We can rewrite the
equation $\Delta y = m\cdot\Delta x$ as a ratio
\[
m = \dfrac{\Delta y}{\Delta x} = 
		\mbox{the rate of change of $y$ with respect 
			to $x$}.
\]
For these reasons we give $m$ its third interpretation as a {\bf rate
  of change}.

\begin{center}
\begin{picture}(320,60)
	\begin{thicklines}
	\put(0,0){\framebox(320,60)
		{\shortstack{\bf For a linear function satisfying
		\boldmath $\Delta y = m\cdot\Delta x$, \\
		\bf the coefficient \boldmath $m$ is \\ [2pt]
		\bf rate of change, slope, and multiplier.}}}
	\end{thicklines}
\end{picture}
\end{center}

We already use $y'$ to denote the rate of change of $y$, so we can now
write $m = y'$ when $y$ is a linear function of $x$.  In that case we
can also write
\[
\Delta y = y' \cdot \Delta x.
\]
This expression should recall a pattern very familiar to you.  (If
not, change $y$ to $S$ and $x$ to $t$!)  It is the fundamental formula
we have been using to calculate future values of $S$, $I$, and $R$.
We can approach the relation between $y$ and $x$ the same way.  That
is, if $y_0$ is an ``initial value'' of $y$, when $x = x_0$, then {\em
  any\/} value of $y$ can be calculated from
\[
y = y_0 + y' \cdot \Delta x \quad \mbox{or} \quad
		y = y_0 + m \cdot \Delta x.
\]
\label{multiplier-end}

\medskip
\noindent
{\bf Units}.\index{units}  Suppose $x$ and $y$ are quantities that are 
%
\mar{If $x$ and $y$ have units, so does $m$}
%
measured in specific units.  If $y$ is a linear function of $x$, with
$\Delta y = m\cdot\Delta x$, then $m$ must have units too.  Since $m$
is the multiplier that ``converts'' $x$ into $y$, the units for $m$
must be chosen so they will convert $x$'s units into $y$'s units.  In
other words,
\[
\mbox{units for $y$} = \mbox{units for $m$} \times 
		\mbox{units for $x$}.
\]
This implies
\[
\mbox{units for $m$} 
    = \dfrac{\mbox{units for $y$}}{\mbox{units for $x$}}.
\]

For example, the multiplier in the Mississippi River problem converts
years to miles, so it must have units of miles per year.  The rate
equation $R' = bI$ in the $S$-$I$-$R$ model is a more subtle example.
It says that $R'$ is a linear function of $I$.  Since $R'$ is measured
in persons per day and $I$ is measured in persons, we must have
\[
\mbox{units for $b$} 
   = \dfrac{\text{units for $R'$}}{\text{units for $I$}}
   = \dfrac{
         \dfrac{\text{persons}}{\text{days}}} {\text{persons}} 
\]


\medskip
\noindent
{\bf Formulas for linear functions}.\index{linear function!formulas}
The expression $\Delta y = m\cdot\Delta x$ declares that $y$ is a
linear function of $x$, but it doesn't quite tell us what $y$ itself
{\em looks like\/} directly in terms of $x$.  In fact, there are
several equivalent ways to write the relation $y = f(x)$ in a formula,
depending on what information we are given about the function.

\medskip
\noindent
\begin{minipage}[b]{3.65in}
{$\bullet$ \bf The initial-value form}.\index{initial-value form} Here
is a very common situation: we know the value of $y$ at an ``initial''
point---let's say $y_0 = f(x_0)$---and we know the rate of
change---let's say it is $m$.  Then the graph is the straight line of
slope $m$ that passes through the point $(x_0, y_0)$.  The formula for
$f$ is
\end{minipage}
\begin{minipage}[b]{1.3in}
\setlength{\unitlength}{.8pt}
\begin{picture}(150,100)(-20,10)   %(60,17)
%initial value form
	\put(0,25){\vector(1,0){160}}
	\put(35,0){\vector(0,1){115}}
	\put(150,14){\footnotesize $x$}
	\put(40,105){\footnotesize $y$}
	\begin{thicklines}
		\put(5,20){\line(3,2){135}}
	\end{thicklines}
	\put(77,68){\circle*{4}}
	\put(119,96){\circle*{4}}
	\put(79,74){\makebox[0pt][r]
		{\footnotesize $(x_0, y_0)$}}
	\put(121,102){\makebox[0pt][r]
		{\footnotesize $(x, y)$}}
	\put(77,68){\line(1,0){42}}
	\put(119,68){\line(0,1){28}}
	\put(100,58){\makebox[0pt]
		{\footnotesize $\Delta x$}}
	\put(122,78){\makebox[0pt][l]
		{\footnotesize $\Delta y$}}
\end{picture}
\end{minipage} 

\[
y = y_0 + \Delta y = y_0 + m\cdot\Delta x = 
		y_0 + m(x - x_0) = f(x).
\]

\medskip
\noindent
What you should note particularly about this formula is that it
expresses $y$ in terms of the initial data $x_0$, $y_0$, and $m$---as
well as $x$.  Since that data consists of a point $(x_0, y_0)$ and a
slope $m$, the initial-value formula is also referred to as the {\bf
  point-slope form}\index{point-slope form} of the equation of a line.
It may be more familiar to you with that name.

\medskip
\noindent
\begin{minipage}[b]{1.3in}
\setlength{\unitlength}{.8pt}
\begin{picture}(150,100)(60,13)     
%interpolation form
	\put(0,25){\vector(1,0){160}}
	\put(35,0){\vector(0,1){115}}
	\put(150,14){\footnotesize $x$}
	\put(40,105){\footnotesize $y$}
	\begin{thicklines}
		\put(5,20){\line(3,2){135}}
	\end{thicklines}
	\put(50,50){\circle*{4}}
	\put(80,70){\circle*{4}}
	\put(119,96){\circle*{4}}
	\put(50,46){\makebox(0,0)[tl]
		{\footnotesize $(x_1, y_1)$}}
	\put(82,76){\makebox[0pt][r]
		{\footnotesize $(x, y)$}}
	\put(119,92){\makebox(0,0)[tl]
		{\footnotesize $(x_2, y_2)$}}
\end{picture}
\end{minipage} \hfill
\begin{minipage}[b]{3.65in}
{$\bullet$ \bf The interpolation form}.\index{interpolation form} This
time we are given the value of $y$ at {\em two\/} points---let's say
$y_1 = f(x_1)$ and $y_2 = f(x_2)$.  The graph is the line that passes
through $(x_1, y_1)$ and $(x_2, y_2)$, and its slope is therefore
\[
m = \dfrac{y_2 - y_1}{x_2 - x_1}.
\]
\end{minipage}

\medskip

\noindent
Now that we know the slope of the graph we can use the point-slope
form (taking $(x_1, y_1)$ as the ``point'', for example) to get the
equation.  We have
\[
y = y_1 + m(x - x_1) = y_1 + 
		\dfrac{y_2 - y_1}{x_2 - x_1}(x - x_1) = f(x).
\]
Notice how, once again, $y$ is expressed in terms of the initial
data---which consists of the two points $(x_1, y_1)$ and $(x_2, y_2)$.

The process of finding values of a quantity between two given values
is called {\bf interpolation}.\index{interpolation} Since our new
expression does precisely that, it is called the interpolation
formula.  (Of course, it also finds values outside the given
interval.)  Since the initial data is a pair of points, the
interpolation formula is also called the {\bf two-point formula} for
the equation of a line.

\medskip
\noindent
\begin{minipage}[b]{1.3in}
\setlength{\unitlength}{.8pt}
\begin{picture}(150,100)(60,-5)  %(-20,-14)     
%slope-intercept form
	\put(0,25){\vector(1,0){160}}
	\put(35,0){\vector(0,1){115}}
	\put(150,14){\footnotesize $x$}
	\put(40,105){\footnotesize $y$}
	\begin{thicklines}
		\put(5,20){\line(3,2){135}}
	\end{thicklines}
	\put(35,40){\circle*{4}}
	\put(95,80){\circle*{4}}
	\put(32,41){\makebox[0pt][r]
		{\footnotesize $b$}}
	\put(97,86){\makebox[0pt][r]
		{\footnotesize $(x, y)$}}
	\put(35,40){\line(1,0){60}}
	\put(95,40){\line(0,1){40}}
	\put(68,30){\makebox[0pt]
		{\footnotesize $\Delta x$}}
	\put(98,56){\makebox[0pt][l]
		{\footnotesize $\Delta y$}}
\end{picture}
\end{minipage} \hfill
\begin{minipage}[b]{3.65in}
{$\bullet$ \bf The slope-intercept form}.\index{slope-intercept form}
This is a special case of the initial-value form that occurs when the
initial $x_0 = 0$.  Then the point $(x_0, y_0)$ lies on the $y$-axis,
and it is frequently written in the alternate form $(0, b)$.  The
number $b$ is called the {\bf \boldmath
  $y$-intercept}.\index{$y$-intercept} The equation is
	\vspace{-5pt}
\[
y = mx + b = f(x).
\]
\end{minipage}

\medskip
\noindent
In the past you may have thought of this as {\em the\/} formula for a
linear function, but for us it is only one of several.  You will find
that we will use the other forms more often.


\subsection{Functions of Several Variables}
\index{function!of several variables}

{\bf Language and notation}.  Many functions depend on more than one
variable.  For example, sunrise depends on the day of the year but it
also depends on the latitude 
%
\mar{A function can have several input variables} 
%
(position north or south of the equator) of the observer.  Likewise,
the crop yield from an acre of land depends on the amount of
fertilizer used, but it also depends on the amount of rainfall, on the
composition of the soil, on the amount of weeding done---to mention
just a few of the other variables that a farmer has to contend with.

The rate equations in the $S$-$I$-$R$ model also provide examples of
functions with more than one input variable.  The equation
\[
I' = .00001\,SI - I/14
\]
says that we need to specify both $S$ and $I$ to find $I'$.  We can
say that
\[
F(S, I) = .00001\,SI - I/14
\]
is a function whose input is the {\bf ordered pair}\index{ordered
  pair} of variables $(S, I)$.  In this case $F$ is given by an
algebraic formula.  While many other functions of several variables
also have formulas---and they are extremely useful---not all functions
do.  The sunrise function, for example, is given by a two-way table
(see page~\pageref{sunrise table}) that shows the time of sunrise for
different days of the year and different latitudes.

As a technical matter it is important to note that the input variables
$S$ and $I$ of the function $F(S, I)$ above appear in a particular
{\em order}, and that order is part of the definition of the function.
For example, $F(1, 0) = 0$, but $F(0, 1) = -1/14$.  (Do you see why?
Work out the calculations yourself.)

\medskip
\noindent
{\bf Parameters}.\index{parameter} Suppose we rewrite the rate
equation for $I'$, replacing .00001 and 1/14 with the general values
$a$ and $b$:
\[
I' = aSI - bI.
\]
This makes it clear that $I'$ depends on $a$ and $b$, too.  But note
that $a$ and $b$ are not variables in quite the same way that $S$ and
$I$ are.  For example, $a$ and $b$ will vary if we switch from one
disease to another or from one population to another.  However, they
will stay fixed while we consider a particular disease in a particular
population. By contrast, $S$ and $I$ will {\em always\/} be treated as
variables.  We call a quantity like $a$ or $b$ a {\bf parameter}.

To emphasize that $I'$ depends on the parameters as well as $S$ and
$I$, we can write $I'$ as the output of a new function
\[
I'= I'(S, I, a, b) = aSI - bI
\]
whose input is the set of {\em four\/} variables 
%
\mar{Some functions\\ depend on parameters}
%
$(S, I, a, b)$, in that order.  The variables $S$, $I$, and $R$ must
also depend on the parameters, too, and not just on $t$.  Thus, we
should write $S(t, a, b)$, for example, instead of simply $S(t)$.  We
implicitly used the fact that $S$, $I$, and $R$ depend on $a$ and $b$
when we discovered there was a threshold for an epidemic
(page~\pageref{`threshold'}).  In exercise~22 of section~1
(page~\pageref{`setthresh'}), you made the relation explicit.  In that
problem you show $I$ will simply decrease over time (i.e., there will
be no ``burst'' of infection) if
\[
S < \dfrac{b}{a}.
\]

There are even more parameters lurking in the $S$-$I$-$R$ problem.  To
uncover them, recall that we needed {\em two\/} pieces of information
to estimate $S$, $I$, and $R$ over time:
\smallskip

1)\quad the rate equations; 

2)\quad the initial values $S_0$, $I_0$, and $R_0$.

\smallskip
\noindent
We used $S_0 = 45400$, $I_0 = 2100$, and $R_0 = 2500$ in the text, but
if we had started with other values then $S$, $I$, and $R$ would have
ended up being different functions of $t$.  Thus, we should really
write
\[
S = S(t, a, b, S_0, I_0, R_0)
\]
to tell a more complete story about the inputs that determine the
output $S$.  Most of the time, though, we do {\em not\/} want to draw
attention to the parameters; we usually write just $S(t)$.

\medskip
\noindent
{\bf Further possibilities}.  Steps I, II, and III on
page~\pageref{`flowchart'} are also functions, because they have
well-defined input and output.  They are unlike the other examples we
have discussed up to this point because they have more than one output
variable.  You should see, though, that there is nothing more
difficult going on here.

In our study of the $S$-$I$-$R$ model it was natural not to separate
functions that have one input variable from those that have several.
This is the pattern we shall follow in the rest of the course.  In
particular, we will want to deal with parameters, and we will want to
understand how the quantities we are studying depend on those
parameters.


\subsection{The Beginnings of Calculus}

While functions, graphs, and computers are part of the general fabric
of mathematics, we can also abstract from the $S$-$I$-$R$ model some
important aspects of the calculus itself.  The first of these is the
idea of a {\bf rate of change}.\index{rate of change} In this chapter
we just assumed the idea was intuitively clear.  However, there are
some important questions not yet answered; for example, how do you
deal with a quantity whose rate of change is itself always changing?
These questions, which lead to the fundamental idea of a {\bf
  derivative},\index{derivative} are taken up in chapter 3.

Rate equations---more\index{rate equation} commonly called {\bf
  differential equations}---lie at the very heart of
calculus.\index{differential equation} We will have much more to say
about them, because many processes in the physical, biological, and
social realms can be modelled by rate equations.  In our analysis of
the $S$-$I$-$R$ model, we used rate equations to estimate future
values by assuming that rates stay fixed 
%
\mar{How the next three chapters are connected} 
%
for a whole day at a time.  The discussion called ``there and back
again'' on page~\pageref{`thereback'} points up the shortcomings of
this assumption.  In chapter 2 we will develop a procedure, called
Euler's method, to address this problem.  In chapter 4 we will return
to differential equations in a general way, equipped with Euler's
method and the concept of the derivative.

\subsection{Exercises}

\subsubsection{Functions and Graphs}

\vspace{-10pt}

\num Sketch the graph of each of the following functions.  Label each
axis, and mark a scale of units on it.  For each line that you draw,
indicate

	i) its slope;

	ii) its $y$-intercept;

	iii) its $x$-intercept (where it crosses the
        $x$-axis).\index{$x$-intercept}

\sub $y =  -\tfrac{1}{2} x + 3$ \hfill c) $5x + 3y = 12$
\hfill \mbox{}

\sub  $y = (2 x - 7)/3$  

\num Graph the following functions.  Put labels and scales on the
axes.

\sub  $V = .3Z -1$; \hfill b)  $W = 600 - P^2$. \hfill \mbox{}

\num Sketch the graph of each of the following functions.  Put labels
and scales on the axes.  For each graph that you draw, indicate

	i) its $y$-intercept;
	
	ii) its $x$-intercept(s).
	
\noindent
For part (d) you will need the {\bf quadratic formula}
\index{quadratic formula}
\[
x = \frac{- b \pm \sqrt{b^2 - 4ac}}{2a}
\] 
for the roots of the {\bf quadratic equation} $ax^2 + bx + c = 0.$

\medskip
\noindent
\begin{minipage}{2in}
a) $y = x^2 $ 	

b) $y = x^2 + 1$
\end{minipage}
\hfill
\begin{minipage}{2in} 
c) $y = (x+1)^2$ 

d) $y = 3x^2 +x -1$
\end{minipage}
\hfill
\mbox{}

\bigskip
\noindent
The next four questions refer to these functions: 
\begin{alignat*}{2}
 c(x, y) &=  17 & &\mbox{a constant function} \\
    j(z) &= z  &  &\mbox{the identity function} \\
    r(u) &= 1/u & & \mbox{the reciprocal function} \\
    D(p,q) &= p - q & &\mbox{the difference function}\\
    s(y) &= y^2 & &\mbox{the squaring function}\\ 
   Q(v) &= \dfrac{2v+1}{3v-6} & &\mbox{a rational function}\\
   H(x) &= \begin{cases}
	5 & \text{if $x < 0$} \\
	x^2 + 2 & \text{if $0 \leq x < 6$} \\
	29 - x & \text{if $6 \leq x$}
	   \end{cases}  & & \\
   T(x, y) &= r(x) + Q(y) & &
\end{alignat*}


\num  Determine the following values: 
\begin{center}
	\begin{tabular}{cccc}
$c(5,-3)$    & $j(17)$   & $c(a,b)$    & $j(u^2+1)$ \\[.05in]
$j(c(3,-5))$ & $s(1.1)$  & $r(1/17)$     & $Q(0)$ \\[.05in]
$Q(2)$       & $Q(3/7)$  & $D(5,-3)$ & $D(-3,5)$ \\[.05in]
$H(1)$       &  $H(7)$   &  $H(4)$ & $H(H(H(-3)))$ \\[.05in]
$r(s(-4))$   & $r(Q(3))$ & $Q(r(3))$ & $T(3, 7)$ 
	\end{tabular}
\end{center}

\num True or false.  Give reasons for your answers: if you say true,
explain why; if you say false, give an example that shows why it is
false.

\sub For every non-zero number $x$, $r(r(x)) = j(x)$.

\sub If $a > 1$, then $s(a) > 1$.

\sub If $a > b$, then $s(a) > s(b)$.

\sub For all real numbers $a$ and $b$, $s(a+b) = s(a) + s(b)$.

\sub  For all real numbers $a$, $b$, and $c$, $D(D(a,b),c)) = D(a,D(b,c))$.

\num  Find all numbers $x$ for which $Q(x) = r(Q(x))$.

\num The {\bf natural
  domain}\index{domain}\index{domain!natural}\index{natural domain} of
a function $f$ is the largest possible set of real numbers $x$ for
which $f(x)$ is defined.  For example, the natural domain of $r(x) =
1/x$ is the set of all non-zero real numbers.

\sub  Find the natural domains of $Q$ and $H$.  

\sub  Find the natural domains of $P(z) = Q(r(z))$; 
$R(v) = r(Q(v))$.


\sub What is the natural domain of the function $W(t) =
\sqrt{\dfrac{1 - t^2}{t^2 - 4}}$?


\subsubsection{Computer Graphing}
\index{computer graphing}\index{graph!computer}

The purpose of these exercises is to give you some experience using a
``graphing package'' on a computer.  This is a program that will draw
the graph of a function $y = f(x)$ whose formula you know.  You must
type in the formula, using the following symbols to represent the
basic arithmetic operations:

\begin{center}
	\begin{tabular}{ccc}
	to indicate & & type \\ \hline
	addition & \rule{0pt}{15pt} & {\tt +} \\
	subtraction & & {\tt -} \\
	multiplication & & {\tt *} \\
	division & & {\tt /} \\
	an exponent & & \verb+^+
	\end{tabular}
\end{center}

\noindent The caret `` \verb+^+ '' appears above the ``{\tt 6}'' on a
keyboard (Shift-6).

\noindent  Here is an example:
\begin{center}
	\begin{tabular}{ccc}
	to enter: & \rule{24pt}{0pt} & type: \\ 
	$\dfrac{7 x^5 - 9 x^2}{x^3 + 1}$
		& \rule{0pt}{20pt} & 
		\verb!(7*x^5 - 9*x^2)/(x^3 + 1)!
	\end{tabular}
\end{center}

\noindent The parentheses you see here are important.  If you do not
include them, the computer will interpret your entry as
\[
7 x^5 - \dfrac{9 x^2}{x^3} + 1 = 
		7 x^5 - \dfrac{9}{x} + 1 \neq 
		\dfrac{7 x^5 - 9 x^2}{x^3 + 1}.
\]
In some graphing packages, you do not need to use {\tt *} to indicate
a multiplication.  If this is true for the package you use, then you
can enter the fractional expression above in the somewhat simpler form

\begin{center}
		\verb!(7x^5 - 9x^2)/(x^3 + 1).!
\end{center}
To do the following exercises, follow the specific instructions for
the graphing package you are using.

\num  Graph the function $f(x) = .6 x + 2$ on the interval $-4 \leq x \leq 4$.  

\sub  What is the $y$-intercept of this graph?  What is the $x$-intercept?

\sub Read from the graph the value of $f(x)$ when $x = -1$ and when $x
= 2$.  What is the difference between these $y$ values?  What is the
difference between the $x$ values?  According to these differences,
what is the slope of the graph?  According to the {\em formula}, what
is the slope?

\num Graph the function $f(x) = 1 - 2 x^2$ on the interval $-1 \leq x
\leq 1$.\label{roots}

\sub What is the $y$-intercept of this graph?  The graph has two
$x$-intercepts; use algebra to find them.

\smallskip

You can also find an $x$-intercept using the computer.  The idea is to
{\bf magnify}\index{magnify} the graph near the intercept until you
can determine as many decimal places in the $x$ coordinate as you
want.  For a start, graph the function on the interval $0\leq
x\leq1$. You should be able to see that the graph on your computer
monitor crosses the $x$-axis somewhere around .7.  Regraph $f(x)$ on
the interval $.6 \leq x \leq .8$.  You should then be able to
determine that the $x$-intercept lies between .70 and .71.  This means
$x =$ .7\ldots ; that is, you know the location of the $x$-intercept
to one decimal place of accuracy.

\sub Regraph $f(x)$ on the interval $.70 \leq x \leq .71$ to get two
decimal places of accuracy in the location of the $x$-intercept.
Continue this process until you have at least 7 places of accuracy.
What is the $x$-intercept?

\bigskip

\noindent
{\bf The circular functions} Graphing packages ``know'' the familiar
functions of trigonometry.  Trigonometric functions are qualitatively
different from the functions in the preceding problems.  Those
functions are defined by algebraic formulas, so they are called {\bf
  algebraic functions}. \index{function!algebraic}\index{algebraic
  function} The trigonometric functions are defined by explicit
``recipes," but {\em not\/} by algebraic formulas; they are called
{\bf transcendental functions}. \index{transcendental
  function}\index{function!transcendental} For calculus, we usually
use the definition of the trigonometric functions \index{trigonometric
  functions}\index{function!trigonometric} as {\bf circular
  functions}. \index{circular functions}\index{function!circular} This
definition begins with a unit circle centered at the origin.  Given
the input number $t$, locate a point $P$ on the circle by tracing an
arc of length $t$ along the circle from the point $(1,0)$.  If $t$ is
positive, trace the arc counterclockwise; if $t$ is negative, trace it
clockwise.  Because the circle has radius $1$, the arc of length $t$
subtends a central angle of {\bf radian} measure $t$.\index{radian}

\special{psfile=../ps/calc1/trig1.ps}

The circular (or trigonometric) functions $\cos t$ and $\sin t$ are
defined as the coordinates of the point $P$,
\[
P = (\cos{t}, \sin{t}).
\]  
The other trigonometric functions are defined in terms of the sine and
cosine:
\begin{alignat*}{2}
\tan{t} &= \sin{t}/\cos{t}, & \qquad \sec{t} &= 1/\cos{t}, \\
\cot{t} &= \cos{t}/\sin{t}, &        \csc{t} &= 1/\sin{t}.
\end{alignat*}


\special{psfile=../ps/calc1/trig2.ps}

Notice that when $t$ is a positive acute angle, the circle definition
agrees with the right triangle definitions of the sine and cosine:
\[
\sin t = \frac{\mbox{opposite}}{\mbox{hypotenuse}} \qquad \mbox{and} \qquad 
\cos t = \frac{\mbox{adjacent}}{\mbox{hypotenuse}}.
\]
However, the circle definitions of the sine and cosine have the
important advantage that they produce functions whose domains are the
set of {\em all\/} real numbers. (What are the domains of the tangent,
secant, cotangent and cosecant functions?)

\enlargethispage*{10pt}
In calculus, angles are also always measured in radians.  To convert
between radians and degrees, notice that the circumference of a unit
circle is $2\pi$, so the radian measure of a semi-circular arc is half
of this, and thus we have
\[
\pi \  \mbox{radians} \; = 180 \  \mbox{degrees}.
\]
As the course progresses, you 
%
\mar{Simplicity determines the choice of radians for measuring angles}
%
will see why radians are used rather than degrees (or mils or any
other unit for measuring angles)---it turns out that the formulas
important to calculus take their simplest form when angles are
expressed in radians.

Graphing packages ``know" the trigonometric functions in exactly this
form: circular functions with the input variable given in radians.
You might wonder, though, how a computer or calculator ``knows'' that
$\sin(1) = .017452406\ldots$.  It certainly isn't drawing a very
accurate circle somewhere and measuring the $y$ coordinate of some
point.  While the circular function approach is 
%
\mar{How do we get\\ values for the\\ circular functions?}  
%
a useful way to think about the trigonometric functions conceptually,
it isn't very helpful if we actually want values of the functions.
One of the achievements of calculus, as you will see later in this
course, is that it provides effective methods for computing values of
functions like the circular functions that aren't given by algebraic
formulas.  \

The following exercises let you review the trigonometric functions and
explore some of the possibilities using computer graphing.

\num Graph the function $f(x) = \sin(x)$ on the interval $-2 \leq x
\leq 10$.

\sub What are the $x$-intercepts of $\sin(x)$ on the interval $-2 \leq
x \leq 10$?  Determine them to two decimal places accuracy.

\sub What is the largest value of $f(x)$ on the interval $-2 \leq x
\leq 10$?  Which value of $x$ makes $f(x)$ largest?  Determine $x$ to
two decimal places accuracy.

\sub Regraph $f(x)$ on the very small interval $-.001 \leq x \leq
.001$.  Describe what you see.  Can you determine the slope of this
graph?

\num Graph the function $f(x) = \cos(x)$ on the interval $0 \leq x
\leq 14$.  On the same screen graph the {\em second\/} function $g(x)
= \cos(2x)$.

\sub How far apart are the $x$-intercepts of $f(x)$?  How far apart
are the $x$-intercepts of $g(x)$?

\sub The graph of $g(x)$ has a pattern that repeats.  How wide is this
pattern?  The graph of $f(x)$ also has a repeating pattern; how wide
is {\em it\/}?

\sub Compare the graphs of $f(x)$ and $g(x)$ to one another.  In
particular, can you say that one of them is a stretched or compressed
version of the other?  Is the compression (or stretching) in the
vertical or the horizontal direction?

\sub Construct a {\em new\/} function $f(x)$ whose graph is the same
shape as the graph of $g(x) = \cos(2x)$, but make the graph of $f(x)$
twice as tall as the graph of $g(x)$.  [A suggestion: either deduce
  what $f(x)$ should be, or make a guess.  Then test your choice on
  the computer.  If your choice doesn't work, think how you might
  modify it, and then test your modifications the same way.]

\num The aim here is to find a solution to the equation $\sin x =
\cos(3x)$.  There is no purely {\em algebraic\/} procedure to solve
this equation.  Because the sine and cosine are not defined by {\em
  algebraic} formulas, this should not be particularly surprising.
(Even for algebraic equations, there are only a few very special cases
for which there are formulas like the quadratic formula. In chapter 5
we will look at a method for solving equations when formulas can't
help us.)

\sub Graph the two functions $f(x) = \sin(x)$ and $g(x) = \cos(3x)$ on
the interval $0 \leq x \leq 1$.

\sub Find a solution of the equation $\sin(x) = \cos(3x)$ that is
accurate to six decimal places.

\sub Find {\em another\/} solution of the equation $\sin(x) =
\cos(3x)$, accurate to four decimal places.  Explain how you found it.


\num Use a graphing program to make a sketch of the graph of each of
the following functions.  In each case, make clear the domain and the
range of the function, where the graph crosses the axes, and where the
function has a maximum or a minimum.
\begin{xalignat*}{2}
&\text{a) } F(w) = (w-1)(w-2)(w-3) &\qquad
	&\text{b) } Q(a) = \dfrac{1}{a^2+5} \\[3pt]
&\text{c) } E(x) = x + \dfrac{1}{x}  &
        &\text{d) } e(x) = x - \dfrac{1}{x}  \\[3pt]
&\text{e) } g(u) = \sqrt{\dfrac{u-1}{u+1}} &
        &\text{f) } M(u) = \dfrac{u^2-2}{u^2+2}  
\end{xalignat*}

\num  Graph on the same screen the following three functions:
\[
f(x) = 2^x, \quad\qquad g(x) = 3^x, \quad\qquad h(x) = 10^x.
\]
Use the interval $-2 \leq x \leq 1.2$.  

\sub  Which function has the largest value when $x = -2$?

\sub  Which is climbing most rapidly when $x = 0$?

\sub Magnify the picture at $x = 0$ by resetting the size of the
interval to $-.0001 \leq x \leq .0001$.  Describe what you see.
Estimate the slopes of the three graphs.


\subsubsection{Proportions, Linear Functions, and Models}
\index{model}\index{linear function}

\vspace{-10pt} 

\num Go back to the three functions given in problem 1.  For each
function, choose an initial value $x_0$ for $x$, find the
corresponding value $y_0$ for $y$, and express the function in the
form $y-y_0=m\cdot(x-x_0)$.

\num You should be able to answer all parts of this problem without
ever finding the equations of the functions involved.

\sub Suppose $y=f(x)$ is a linear function with multiplier $m=3$.  If
$f(2)=-5$, what is $f(2.1)$? $f(2.0013)$?  $f(1.87)$? $f(922)$?

\sub Suppose $y=G(x)$ is a linear function with multiplier $m=-2$.  If
$G(-1)=6$, for what value of $x$ is $G(x)=8$? $G(x)=0$? $G(x)=5$?
$G(x)=491$?

\sub Suppose $y=h(x)$ is a linear function with $h(2)=7$ and $h(6)=9$.  What is $h(2.046)$? $h(2+a)$?

\num In Massachusetts there is a sales tax of 5\%.\index{sales tax}
The tax $T$, in dollars, is proportional to the price $P$ of an
object, also in dollars.  The constant of proportionality is $k = 5\%
= .05$.  Write a formula that expresses the sales tax as a linear
function of the price, and use your formula to compute the tax on a
television set that costs \$289.00 and a toaster that costs \$37.50.

\num Suppose $W = 213 - 17\, Z$.  How does $W$ change when $Z$ changes
from 3 to 7; from 3 to 3.4; from 3 to 3.02?  Let $\Delta Z$ denote a
change in $Z$ and $\Delta W$ the change thereby produced in $W$.  Is
$\Delta W = m\,\Delta Z$ for some constant $m$?  If so, what is $m$?

\numa In the following table, $q$ is a linear function of $p$.  Fill
in the blanks in the table.
\begin{center}
	\begin{tabular}
		{c|c|c|c|c|c|c|c}
	\makebox[18pt]{$p$} & \makebox[18pt]{--3} &
 	\makebox[18pt]{0}   & \makebox[18pt]{}    &
	\makebox[18pt]{7}   & \makebox[18pt]{13}  &
	\makebox[18pt]{}    & \makebox[18pt]{$\pi$} \\ [2pt] 
		\hline 
	\makebox[18pt]{$q$} & \makebox[18pt]{7}   &
	\rule{0pt}{14pt}    & \makebox[18pt]{4}   &
	\makebox[18pt]{1}   & \makebox[18pt]{}    &
	\makebox[18pt]{0}   & \makebox[18pt]{}    
	\end{tabular}
\end{center}

\sub Find a formula to express $\Delta q$ as a function of $\Delta p$,
and another to express $q$ as a function of $p$.

\num {\bf Thermometers}.\index{thermometer} There are two scales in
common use to measure the temperature, called {\bf Fahrenheit degrees}
and the {\bf Celsius degrees}.\index{Fahrenheit}\index{Celsius} Let
$F$ and $C$, respectively, be the temperature on each of these scales.
Each of these quantities is a linear function of the other; the
relation between them in determined by the following table:
\begin{center}
	\begin{tabular}{ccc}
	physical measurement & $C$ & $F$ \\ [1pt] \hline
	freezing point of water & 
		\rule{0pt}{14pt} 0 \rule{0pt}{14pt} & 32 \\
	boiling point of water & 100 & 212 
	\end{tabular}
\end{center}

\sub Which represents a larger change in temperature, a Celsius degree
or a Fahrenheit degree?

\sub How many Fahrenheit degrees does it take to make the temperature
go up one Celsius degree?  How many Celsius degrees does it take to
make it go up one Fahrenheit degree?

\sub What is the multiplier $m$ in the equation $\Delta F = m \cdot
\Delta C$?  What is the multiplier $\mu$ in the equation $\Delta C =
\mu \cdot \Delta F$?  (The symbol $\mu$ is the Greek letter {\em mu}.)
What is the relation between $\mu$ and $m$?

\sub Express $F$ as a linear function of $C$.  Graph this function.
Put scales and labels on the axes.  Indicate clearly the slope of the
graph and its vertical intercept.

\sub Express $C$ as a linear function of $F$ and graph this function.
How are the graphs in parts (d) and (e) related?  Give a clear and
detailed explanation.

\sub Is there any temperature that has the same reading on the two
temperature scales?  What is it?  Does the temperature of the air ever
reach this value?  Where?


\num {\bf The Greenhouse Effect}.\index{greenhouse effect} The
concentration of carbon dioxide\index{carbon dioxide} (CO$_2$) in the
atmosphere is increasing.  The concentration is measured in parts per
million (PPM).  Records kept at the South Pole show an increase of .8
PPM per year during the 1960s.

\sub At that rate, how many years does it take for the concentration
to increase by 5 PPM; by 15 PPM?

\sub At the beginning of 1960 the concentration was about 316 PPM.
What would it be at the beginning of 1970; at the beginning of 1980?

\sub Draw a graph that shows CO$_2$ concentration as a function of the
time since 1960.  Put scales on the axes and label everything clearly.

\sub The {\em actual\/} CO$_2$ concentration at the South Pole was 324
PPM at the beginning of 1970 and 338 PPM at the beginning of 1980.
Plot these values on your graph, and compare them to your calculated
values.

\sub Using the actual concentrations in 1970 and 1980, calculate a new
rate of increase in concentration.  Using that rate, estimate what the
increase in CO$_2$ concentration was between 1970 and 1990.  Estimate
the CO$_2$ concentration at the beginning of 1990.

\sub Using the rate of .8 PPM per year that held during the 1960s,
determine how many years before 1960 there would have been {\em no\/}
carbon dioxide at all in the atmosphere.


\num {\bf Thermal Expansion}.\index{thermal expansion} Measurements
show that the length of a metal bar increases in proportion to the
increase in temperature.  An aluminum bar that is exactly 100 inches
long when the temperature is 40$^\circ$F becomes 100.0052 inches long
when the temperature increases to 80$^\circ$F.

\sub How long is the bar when the temperature is 60$^\circ$F?
100$^\circ$F?

\sub What is the multiplier that connects an increase in length
$\Delta L$ to an increase in temperature $\Delta T$?

\sub Express $\Delta L$ as a linear function of $\Delta T$.

\sub How long will the bar be when $T = 0^\circ$F?

\sub Express $L$ as a linear function of $T$.

\sub What temperature change would make $L$ = 100.01 inches?

\sub For a {\em steel\/} bar that is also 100 inches long when the
temperature is 40$^\circ$ F, the relation between $\Delta L$ and
$\Delta T$ is $\Delta L = .00067 \, \Delta T$.  Which expands more
when the temperature is increased; aluminum or steel?

\sub How long will this steel bar be when $T = 80^\circ$F?

\num {\bf Falling Bodies}\label{`falling'}.  In the simplest model of
the motion of a falling body, the velocity increases in proportion to
the increase in the time that the body has been falling.  If the
velocity is given in feet per second, measurements show the constant
of proportionality is approximately~32.

\sub A ball is falling at a velocity of 40 feet/sec after 1 second.
How fast is it falling after 3 seconds?

\sub Express the change in the ball's velocity $\Delta v$ as a linear
function of the change in time $\Delta t$.

\sub  Express $v$ as a linear function of $t$.

\noindent
The model can be expanded to keep track of the {\em distance\/} that
the body has fallen.  If the distance $d$ is measured in feet, the
units of $d'$ are feet per second; in fact, $d'=v$.  So the model
describing the motion of the body is given by the rate equations
\begin{align*}
	d' &= v  \quad \mbox{feet per second};\\
	v' &= 32 \quad \mbox{feet per second per second}.
\end{align*}

\sub At what rate is the distance increasing after 1 second? After 2
seconds?  After 3 seconds?

\sub  Is $d$ a linear function of $t$?  Explain your answer. 

\bigskip

In many cases, the rate of change of a variable quantity is
proportional to the quantity itself.  Consider a human population as
an example.  If a city of 100,000 is increasing at the rate of 1500
persons per year, we would expect a similar city of 200,000 to be
increasing at the rate of 3000 persons per year.  That is, if $P$ is
the population at time $t$, then the {\bf net growth
  rate}\index{growth rate!net} $P'$ is proportional to $P$:
\[
P' = k\, P.
\]
In the case of the two cities, we have
\[
P' = 1500 = k\, P = k \times 100000 \quad 
		\mbox{so} \quad k =
		\dfrac{1500}{100000} = .015.
\]

\num In the equation $P' = k\, P$, above, explain why the units for
$k$ are
\[
\dfrac{\text{persons per year}}{\text{person}}.
\]
The number $k$ is called the {\bf per capita growth
  rate}.\index{growth rate!per capita} (``Per capita'' means ``per
person''---``per {\em head\/}'', literally.)

\num {\bf Poland and Afghanistan}.\index{Poland}\index{Afghanistan}\label{Poland} In
1985 the per capita growth rate in Poland was 9 persons per year per
thousand persons.  (That is, $k = 9/1000 = .009$.)  In Afghanistan it
was 21.6 persons per year per thousand.

\sub Let $P$ denote the population of Poland and $A$ the population of
Afghanistan.  Write the equations that govern the growth rates of
these populations.

\sub In 1985 the population of Poland was estimated to be 37.5 million
persons, that of Afghanistan 15 million.  What are the net growth
rates $P'$ and $A'$ (as distinct from the {\em per capita\/} growth
rates)?  Comment on the following assertion: When comparing two
countries, the one with the larger per capita growth rate will have
the larger net growth rate.

\sub On the average, how long did it take the population to increase
by one person in Poland in 1985?  What was the corresponding time
interval in Afghanistan?


\numa {\bf Bacterial Growth}.\index{growth!bacterial} A colony of
bacteria on a culture medium grows at a rate proportional to the
present size of the colony.  When the colony weighed 32 grams it was
growing at the rate of 0.79 grams per hour.  Write an equation that
links the growth rate to the size of the population.

\sub What is $\Delta P$ if $\Delta t$ = 1~minute?  Estimate how long
it would take to make $\Delta P$ = .5~grams.

\num {\bf Radioactivity}.\index{radioactivity} In radioactive decay,
radium slowly changes into lead.  If one sample of radium is twice the
size of a second lump, then the larger sample will produce twice as
much lead as the second in any given time.  In other words, the rate
of decay is proportional to the amount of radium present.
Measurements show that 1 gram of radium decays into lead at the rate
of 1/2337 grams per year.  Write an equation that links the decay rate
to the size of the radium sample.  How does your equation indicate
that the process involves {\em decay\/} rather than {\em growth}?

\num {\bf Cooling}.\index{cooling}\label{`cooling'} Suppose a cup of
hot coffee is brought into a room at 70$^{\circ}$F.  It will cool off,
and it will cool off {\em faster\/} when the temperature difference
between the coffee and the room is greater.  The simplest assumption
we can make is that the rate of cooling is proportional to this
temperature difference (this is called Newton's law of
cooling).\index{Newton, I.}\index{Newton's law of cooling} Let $C$
denote the temperature of the coffee, in $^{\circ}$F, and $C'$ the
rate at which it is cooling, in $^{\circ}$F per minute.  The new
element here is that $C'$ is proportional, not to $C$, but to the {\em
  difference\/} between $C$ and the room temperature of 70$^{\circ}$F.

\sub Write an equation that relates $C'$ and $C$.  It will contain a
proportionality constant $k$.  How did you indicate that the coffee is
{\em cooling\/} and not {\em heating up\/}?

\sub When the coffee is at 180$^{\circ}$F it is cooling at the rate of
9$^{\circ}$F per minute.  What is $k$\/?

\sub At what rate is the coffee cooling when its temperature is
120$^{\circ}$F?

\sub Estimate how long it takes the temperature to fall from
180$^{\circ}$F to 120$^{\circ}$F.  Then make a better estimate, and
explain why it is better.

\newpage
%%%%%%%%%%%%%%%%%%%%%%%%%%%%%%%%%%%%%%%%%%%%%%%%%%%%%%%%%%%%%%%%%%%%%%%
%                                                                     %
%                             Section 3                               %
%                                                                     % 
%%%%%%%%%%%%%%%%%%%%%%%%%%%%%%%%%%%%%%%%%%%%%%%%%%%%%%%%%%%%%%%%%%%%%%%

\section{Using a Program}
\index{program}

\subsection{Computers}
\index{computer}

A computer changes the way we can use calculus as a tool, and it
vastly enlarges the range of questions that we can tackle.  No longer
need we back away from a problem that involves a lot of computations.
There are two aspects to the power of a computer.  First, it is fast.
It can do a million additions in the time it takes us to do one.
Second, it can be programmed.  By arranging computations into a
loop---as we did on page~\pageref{`loop'}---we can construct a program
with only a few instructions that will carry out millions of
repetitive calculations.

The purpose of this section is to give you practice using a computer
program that estimates values of $S$, $I$, and $R$ in the epidemic
model.  As you will see, it carries out the three rounds of
calculations you have already done by hand.  It also contains a loop
that will allow you to do a hundred, or a million, rounds of
calculations with no extra effort.


\subsubsection{The Program SIR}

The program on the following page calculates values of $S$, $I$,
and~$R$.  It is a set of instructions---sometimes called {\bf
  code}---that is designed to be read by you and by a
computer.\index{code} These instructions mirror the operations we
performed by hand to generate the table on page~\pageref{`threeday'}.
\mar{Read a program line by line from the top} The code here is
similar to what it would be in most programming languages.  The line
numbers, however, are not part of the program; they are there to help
us refer to the lines.  A computer reads the code one line at a time,
starting at the top.  Each line is a complete instruction which causes
the computer to do something.  The purpose of nearly every instruction
in this program is to assign a numerical value to a symbol.  Watch for
this as we go down the lines of code.

The first line, {\tt t = 0}, 
%
\mar{Lines 1--5}
%
is the instruction ``Give $t$ the value 0.''  The next four lines are
similar.  Notice, in the fifth line, how $\Delta t$ is typed out as
{\tt deltat}.  It is a common practice for the name of a variable to
be several letters long.  A few lines later $S'$ is typed out as {\tt
  Sprime}, for instance.  The instruction on the sixth line \mar{Line
  6} is the first that does not assign a value to a symbol.  Instead,
it causes the computer to print the following on the computer monitor
screen:
\begin{center}
{\tt 0} \qquad {\tt 45400} \qquad {\tt 2100} \qquad {\tt 2500}
\end{center}

\newpage

\begin{center}
{\bf Program: SIR}\index{program!SIR}\label{prog}
\end{center}

\vspace{-10pt}

\noindent
\begin{minipage}[t]{.2in}
{\footnotesize\sf 
\vspace{-15pt}
\mbox{}\\
\hphantom{1}1\\ 
\hphantom{1}2\\ 
\hphantom{1}3\\ 
\hphantom{1}4\\ 
\hphantom{1}5\\ 
\hphantom{1}6\\ 
\hphantom{1}7\\ 
\hphantom{1}8\\ 
\hphantom{1}9\\ 
10\\ 11\\ 12\\ 13\\ 14\\ 15\\ 16\\ 17\\ 18 \\19
}
\end{minipage}
\begin{minipage}[t]{3.7in}
\begin{verbatim}
  t = 0
  S = 45400
  I = 2100
  R = 2500
  deltat = 1
  PRINT t, S, I, R
  FOR k = 1 TO 3
          Sprime = -.00001 * S * I
          Iprime = .00001 * S * I - I / 14
          Rprime = I / 14
          deltaS = Sprime * deltat
          deltaI = Iprime * deltat
          deltaR = Rprime * deltat
          t = t + deltat
          S = S + deltaS
          I = I + deltaI
          R = R + deltaR
          PRINT t, S, I, R
  NEXT k
\end{verbatim}
\end{minipage}
\begin{minipage}[t]{1in}
\mbox{}
\vspace{1.3in}

$
\left.
\rule{0pt}{20pt}
\right\} \mbox{Step I}
$

\end{minipage}
\bigskip

\smallskip

Skip over the line that says {\tt FOR k = 1 TO 3}.  
%
\mar{Line 7}
%
It will be easier to understand after we've read the rest of the
program.

Look at the first three indented lines.  
%
\mar{Lines 8--10}
%
You should recognize them as coded versions of the rate equations
\begin{align*}
S' &= -.00001\,SI, \\
I' &= .00001\,SI - I/14, \\
R' &= I/14,
\end{align*}
for the measles epidemic.  (The program uses {\tt *} to denote
multiplication.)  They are instructions to assign numerical values to
the symbols $S'$, $I'$, and $R'$.  For instance, {\tt Sprime = -.00001
  * S * I} (line 8) says
\begin{center}
Give $S'$ the value $-.00001\,SI$; \\
use the current values of $S$ and $I$ to get $-.00001\,SI$. \\
		\end{center}
Now the computer knows that the current values of $S$ and $I$ are
45400 and 2100, respectively.  (Can you see why?)  So it calculates
the product $-.00001 \times 45400 \times 2100 = -953.4$ and then gives
$S'$ the value $-953.4$.  There is an extra step to calculate the
product.

Notice that the first three indented lines are bracketed together and
labelled ``Step I,'' because they carry out Step I in the flow chart.
%
\mar{Lines 11--13}
%
The next three indented lines carry out Step II in the flow chart.
They assign values to three more symbols---namely $\Delta S$, $\Delta
I$, and $\Delta R$---using the current values of $S'$, $I'$, $R'$ and
$\Delta t$.

The next four indented lines present a puzzle.  
%
\mar{Lines 14--17}
%
They don't make sense if we read them as ordinary mathematics.  For
example, in an expression like \mbox{\texttt{t = t + deltat}}, we would cancel
the {\tt t}'s and conclude {\tt deltat = 0}.  The lines {\em do\/}
make sense when we read them as computer instructions, however.  As a
computer instruction, {\tt t = t + deltat} says
\begin{center}
Make the new value of $t$ equal to the current value of $t + \Delta t$.
\end{center}
(To make this clear, some computer languages express this instruction
in the form {\tt let t = t + deltat}.)  Once again we have an
instruction that assigns a numerical value to a symbol, but this time
the symbol ($t$, in this case) already has a value before the
instruction is carried out.  The instruction gives it a {\em new\/}
value.  (Here the value of $t$ is changed from 0 to 1.)  Likewise, the
instruction {\tt S = S + deltaS} gives $S$ a new value.  What was the
old value, and what is the new?

Compare the three lines of code 
%
\mar{How a program computes new values}
%
that produce new values of $S$, $I$, and $R$ with the original
equations that we used to define Step III back on
page~\pageref{`flowchart'}:
\begin{alignat*}{2}
&\texttt{S = S + deltaS} &\qquad\qquad  \text{new $S$} &= 
           \text{current $S$ + $\Delta S$}, \\
&\texttt{I = I + deltaI} &              \text{new $I$} &= 
           \text{current $I$ + $\Delta I$}, \\
&\texttt{R = R + deltaR} &              \text{new $R$} &= 
           \text{current $R$ + $\Delta R$}. 
\end{alignat*}
The words ``new'' and ``current'' aren't needed in the computer code
because they are automatically understood to be there.  Why?  First of
all, a symbol (like {\tt S}) always has a {\em current\/} value, but
an instruction can give it a {\em new\/} value.  Second, a computer
instruction of the form {\tt A = B} is always understood to mean ``new
{\tt A} = current {\tt B}.''


\aside{Notice that the instructions {\tt A = B} and {\tt B = A} mean
  different things.  The second says ``new {\tt B} = current {\tt
    A}.''  Thus, in {\tt A = B}, {\tt A} is altered to equal {\tt B},
  while in {\tt B = A}, {\tt B} is altered to equal {\tt A}.  To
  emphasize that the symbol on the left is always the one affected,
  some programming languages use a modified equal sign, as in {\tt A
    := B}.  We sometimes read this as ``{\tt A} gets {\tt B}''.}

\newpage

The next line is another {\tt PRINT} statement, 
%
\mar{Line 18}
%
exactly like the one on line 6.  It causes the current values of $t$,
$S$, $I$, and $R$ to be printed on the computer monitor screen.  But
this time what appears is
\begin{center}
{\tt 1} \qquad {\tt 44446.6} \qquad {\tt 2903.4} \qquad {\tt 2650}
\end{center}
The values were changed by the previous four instructions.  It is
important to remember that the computer carries out instructions in
the order they are written.  Had the second {\tt PRINT} statement
appeared right after line 13, say, the old values of $t$, $S$, $I$,
and $R$ would have appeared on the monitor screen a second time.

We will take the last line and line 7 together.  
%
\mar{Lines 7 and 19:\\ the loop}
%
They are the instructions for the loop.  Consider the situation when
we reach the last line.  The variables $t$, $S$, $I$, and $R$ now have
their ``day 1'' values.  To continue, we need an instruction that will
get us back to line 8, because the instructions on lines 8--17 will
convert the current (day 1) values of $t$, $S$, $I$, and $R$ into
their ``day 2'' values.  That's what lines 7 and 19 do.

Here is the meaning of the instruction {\tt FOR k = 1 TO 3} on line 7:
\begin{center}
	Give $k$ the value 1, 
	\mar{\mbox{\tt FOR k = 1 TO 3}}
		and be prepared later to  \\
	give it the value 2 and then the value 3.
\end{center}
The variable $k$ plays the role of a {\bf counter},\index{counter}
telling us how many times we have gone around the loop.  Notice that
$k$ did not appear in our hand calculations.  However, when we said we
had done three rounds of calculations, for example, we were really
saying $k = 3$.

After the computer reads and executes line 7, it carries out all the
instructions from lines 8 to 18, arriving finally at the last line.
The computer then interprets the instruction {\tt NEXT k} on the last
line as follows:
\begin{center}
	Give $k$ the next value 
	\mar{\tt NEXT k}
		that the {\tt FOR} command allows, and \\
	move back to the line immediately after the {\tt FOR}
		command.
\end{center}
After the computer carries out this instruction, $k$ has the value 2
and the computer is set to carry out the instruction on line 8.  It
then executes that instruction, and continues down the program, line
by line, until it reaches line 19 once again.  This sets the value of
$k$ to 3 and moves the computer back to line 8.  Once again it
continues down the program to line 19.  \mar{How the program stops}
This time there is no allowable value that $k$ can be given, so the
program stops.

The {\tt NEXT k} command is different from all the others in the
program.  It is the only one that directs the computer to go to a
different line.  That action causes the program to {\bf loop}.
Because the loop involves all the indented instructions between the
{\tt FOR} statement and the {\tt NEXT} statement, it is called a {\bf
  FOR--NEXT loop}.\index{FOR--NEXT loop} This is just one kind of
loop.  Computer programs can contain other sorts that carry out
different tasks.  In the next chapter we will see how a {\bf
  DO--WHILE} loop is used.



\subsection{Exercises}

\subsubsection{The program SIR}

The object of these exercises is to verify that the program SIR works
the way the text says it does.  Follow the instructions for running a
program on the computer you are using.


\num Run the program to confirm that it reproduces what you have
already calculated by hand (table, page~\pageref{`threeday'}).


\num On a copy of the program, mark the instructions that carry out
the following tasks:

\sub give the input values of $S$, $I$, and $R$;

\sub say that the calculations take us 1 day into the future;

\sub carry out step II (see page~\pageref{`flowchart'});

\sub carry out step III;

\sub give us the output values of $S$, $I$, and $R$;

\sub take us once around the whole loop;

\sub say how many times we go around the loop.


\num Delete all the lines of the program from line 7 onward (or else
type in the first 6 lines).  Will this program run?  What will it do?
Run it and report what you see.  Is this what you expected?


\num Starting with the original SIR program on page~\pageref{prog},
delete lines 7 and 19.  These are the ones that declare the FOR--NEXT
loop.  Will this program run?  What will it do?  Run it and report
what you see.  Is this what you expected?


\num Using the 17-line program you constructed in the previous
question, remove the PRINT statement from the last line and insert it
between what appear as lines 13 and 14 on page~\pageref{prog}.  Will
this program run?  What will it do?  Run it and report what you see.
Is this what you expected?


\num Starting with the original SIR program on page~\pageref{prog},
change line 7 so it reads {\tt FOR k = 26 TO 28}.  Thus, the counter
$k$ takes the values 26, 27, and 28.  Will this program run?  What
will it do?  Run it and report what you see.  Is this what you
expected?



\subsubsection{Programs to practice on}

In this section there are a number of short programs for you to analyze and run.
\begin{center}
\begin{tabular}[t]{ccc}
\begin{minipage}[t]{1.3in}
\centering{\bf Program 1} \\
\begin{verbatim}
  A = 2
  B = 3
  A = B
  PRINT A, B
\end{verbatim}
\end{minipage}
	&
\begin{minipage}[t]{1.3in}
\centering{\bf Program 2} \\
\begin{verbatim}
  A = 2
  B = 3
  B = A
  PRINT A, B
\end{verbatim}
\end{minipage}
	&
\begin{minipage}[t]{1.3in}
\centering{\bf Program 3} \\
\begin{verbatim}
  A = 2
  B = 3
  A = A + B
  B = A + B
  PRINT A, B
\end{verbatim}
\end{minipage}
\end{tabular}
\end{center}

\num When Program 1 runs it will print the values of $A$ and $B$ that
are current when the program stops.  What values will it print?  Type
in this program and run it to verify your answers.

\num What will Program 2 do when it runs? Type in this program and run
it to verify your answers.

\num After each line in Program 3 write the values that $A$ and $B$
have {\em after\/} that line has been carried out.  What values of $A$
and $B$ will it print?  Type in this program and run it to verify your
answers.

\bigskip
\noindent
The next three programs have an element not found in the program SIR.
In each of them, there is a FOR--NEXT loop, and the counter $k$
actually appears in the statements within the loop.

\begin{center}
\begin{tabular}[t]{ccc}
\begin{minipage}[t]{1.5in}
\centering{\bf Program 4} \\
\begin{verbatim}
  FOR k = 1 TO 5
       A = k ^ 3
       PRINT A
  NEXT k
\end{verbatim}
\end{minipage}
	&
\begin{minipage}[t]{1.5in}
\centering{\bf Program 5} \\
\begin{verbatim}
  FOR k = 1 TO 5
       A = k ^ 3
  NEXT k
  PRINT A
\end{verbatim}
\end{minipage}
	&
\begin{minipage}[t]{1.5in}
\centering{\bf Program 6} \\
\begin{verbatim}
  x = 0
  FOR k = 1 TO 5
       x = x + k
       PRINT k, x
  NEXT k
\end{verbatim}
\end{minipage}
\end{tabular}
\end{center}

\num What output does Program 4 produce?  Type in the code and run the
program to confirm your answer.

\num What is the difference between the code in Program 5 and the code
in Program 4?  What is the output of Program 5?  Does it differ from
the output of Program 4?  If so, why?

\num What output does Program 6 produce?  Type in the code and run the
program to confirm your answer.


\begin{center}
\begin{tabular}[t]{ccc}
\begin{minipage}[t]{1.6in}
\centering{\bf Program 7} \\
\begin{verbatim}
  A = 0
  B = 0
  FOR k = 1 TO 5
       A = A + 1
       B = A + B
       PRINT A, B
  NEXT k
\end{verbatim}
\end{minipage}
	&
\begin{minipage}[t]{1.6in}
\centering{\bf Program 8} \\
\begin{verbatim}
  A = 0
  B = 1
  FOR k = 1 TO 5
       A = A + B
       B = A + B
       PRINT A, B
  NEXT k
\end{verbatim}
\end{minipage}
	&
\begin{minipage}[t]{1.5in}
\centering{\bf Program 9} \\
\begin{verbatim}
  A = 0
  B = 1
  FOR k = 1 TO 5
       A = A + B
       B = A + B
  NEXT k
  PRINT A, B
\end{verbatim}
\end{minipage}
\end{tabular}
\end{center}

\num Program 7 prints five lines of output.  What are they?  Type in
the program and run it to confirm your answers.

\num What is the output of Program 8?  Type in the program and run it
to confirm your answers.

\num Describe exactly how the {\em codes\/} for Programs 8 and 9
differ.  How do the {\em outputs\/} differ?


\subsubsection{Analyzing the measles epidemic}

\vspace{-10pt} 

\num Alter the program SIR to have it calculate estimates for $S$,
$I$, and $R$ over the first {\em six\/} days.  Construct a table that
shows those values.\index{estimate}


\num Alter the program to have it estimate the values of $S$, $I$, and
$R$ for the first {\em thirty\/} days.

\sub What are the values of $S$, $I$, and $R$ when $t = 30$?

\sub According to these figures, on what day does the infection peak?
What values do you get for $S$, $I$, and $R$?

\sub How can you reconcile the value you just got for $S$ with the
value we obtained algebraically on page~\pageref{`threshold'}?


\num By adding an appropriate PRINT statement after line 10 you can
also get the program to print values for $S'$, $I'$, and $R'$.  Do
this, and check that you get the values shown in the table on
page~\pageref{`threeday'}.

\num According to these estimates, on what day do the largest number
new infections occur?  How many are there?  Explain where you got your
information.

\num On what day do you estimate that the largest number of recoveries
occurs?  Do you see a connection between this question and 17~(b)?

\num On what day do you estimate the infected population grows most
rapidly?  Declines most rapidly?  What value does $I'$ have on those
days?

\numa Alter the original SIR program so that it will go {\em
  backward\/} in time, with time steps of 1 day.  Specify the changes
you made in the program.  Use this altered program to obtain estimates
for the values of $S$, $I$, and $R$ yesterday.  Compare your estimates
with those in the text (page~\pageref{`threeday'}).

\sub Estimate the values of $S$, $I$, and $R$ {\em three\/} days
before today.


\num According to the $S$-$I$-$R$ model, when did the infection begin?
That is, how many days before today was the estimated value of $I$
approximately 1?


\num {\bf There and back again}.\index{there and back again} Use the
SIR program, modified as necessary, to carry out the calculations
described in exercise 18 on page~\pageref{`therebackprob'}.  Do your
computer results agree with those you obtained earlier?

\num In exercises 22--26 at the end of section~2
(pages~\pageref{`falling'}--\pageref{`cooling'}) you set up rate
equations to model some other systems.  Choose a couple of these and
think of some interesting questions that could be answered using a
suitably modified SIR program.  Make the modifications and report your
results.

%%%%%%%%%%%%%%%%%%%%%%%%%%%%%%%%%%%%%%%%%%%%%%%%%%%%%%%%%%%%%%%%%%%%%%%
%                                                                     %
%                             Section 4                               %
%                                                                     % 
%%%%%%%%%%%%%%%%%%%%%%%%%%%%%%%%%%%%%%%%%%%%%%%%%%%%%%%%%%%%%%%%%%%%%%%

\section{Chapter Summary}

\subsection{The Main Ideas}

\begin{itemize}

\item Natural processes like the spread of disease can often be
  described by {\bf mathematical models}.  Initially, this involves
  identifying {\bf numerical quantities} and relations between them.

\item A relation between quantities often takes the form of a {\bf
  function}.  A function can be described in many different ways; {\bf
  graphs}, {\bf tables}, and {\bf formulas} are among the most common.

\item {\bf Linear functions} make up a special but important class.
  If $y$ is a linear function of $x$, then {\bf \boldmath $\Delta y =
    m\cdot\Delta x$}, for some constant $m$.  The constant $m$ is a
  {\bf multiplier}, {\bf slope}, and {\bf rate of change}.

\item If $y = f(x)$, then we can consider the {\bf rate of change
  \boldmath $y'$} of $y$ with respect to $x$.  A mathematical model
  whose variables are connected by {\bf rate equations} can be
  analyzed to predict how those variables will change.

\item Predicted changes are {\bf estimates} of the form {\bf \boldmath
  $\Delta y = y'\cdot\Delta x$}.

\item The computations that produce estimates from rate equations can
  be put into a {\bf loop}, and they are readily carried out on a {\bf
    computer}.

\item A computer increases the {\bf scope} and {\bf complexity} of the
  problems we can consider.


\end{itemize}

\subsection{Expectations}

\begin{itemize}  

\item You should be able to work with functions given in various
  forms, to find the output for any given input.

\item You should be able to read a graph.  You should also be able to
  construct the graph of a linear function directly, and the graph of
  a more complicated function using a computer graphing package.

\item You should be able to determine the natural domain of a function
  given by a formula.

\item You should be able to express proportional quantities by a
  linear function, and interpret the constant of proportionality as a
  multiplier.

\item Given any two of these quantities for a linear
  function---multiplier, change in input, change in output---you
  should be able to determine the third.

\item You should be able to model a situation in which one variable is
  proportional to its rate of change.

\item Given the value of a quantity that depends on time and given its
  rate of change, you should be able to estimate values of the
  quantity at other times.

\item For a set of quantities determined by rate equations and initial
  conditions, you should be able to estimate how the quantities
  change.

\item Given a set of rate equations, you should be able to determine
  what happens when one of the quantities reaches a maximum or
  minimum, or remains unchanged over time.

\item You should be able to understand how a computer program with a
  FOR--NEXT loop works.

\end{itemize}


\subsection{Chapter Exercises}

\subsubsection{A Model of an Orchard}
\index{orchard}\index{model}

If an apple orchard occupies one acre of land, how many trees should
it contain so as to produce the largest apple crop?  This is an
example of an {\bf optimization}\index{optimization} problem.  The
word {\em optimum\/}\index{optimum} means ``best
possible''---especially, the best under a given set of conditions.
These exercises seek an optimum by analyzing a simple mathematical
model of the orchard.  The model is the function that describes how
the total yield depends on the number of trees.

An immediate impulse is just to plant a lot of trees, on the
principle: more trees, more apples.  But there is a catch: if there
are too many trees in a single acre, they crowd together.  Each tree
then gets less sunlight and nutrients, so it produces fewer apples.
For example, the relation between the {\em yield per tree}, $Y$, and
the {\em number of trees}, $N$, may be like that shown in the graph
drawn on the left, below.

{\small 
\setlength{\unitlength}{.5pt}
\begin{picture}(670,170)(0,-25)

\put(40,0)
{
\begin{picture}(300,100)
	\put(0,0){\vector(1,0){270}}
	\put(0,0){\vector(0,1){120}}
	\put(260,8){\makebox[0pt]{\small $N$}}
	\put(100,-18){\makebox[0pt]{\small 100}}
	\put(200,-18){\makebox[0pt]{\small 200}}
	\put(7,110){\makebox[0pt][l]{\small $Y$}}
	\put(-7,113){\makebox[0pt][r]{\small pounds}}
	\put(-7,94){\makebox[0pt][r]{\small per tree}}
	\put(-6,67){\makebox[0pt][r]{\small 750}}
	\put(260,-18){\makebox[0pt]{\small trees}}
	\begin{thicklines}
		\put(0,70){\line(1,0){20}}
		\bezier{75}(20,70)(40,70)(100,40)
		\bezier{150}(100,40)(180,0)(250,0)
	\end{thicklines}
%	\multiput(20,70)(5,0){4}{\circle*{1}}
%	\multiput(40,70)(3.25,-1.625){43}{\circle*{1}}
\end{picture}
}
\put(370,0)
{
\begin{picture}(300,100)
	\put(0,0){\vector(1,0){270}}
	\put(0,0){\vector(0,1){120}}
	\put(260,8){\makebox[0pt]{\small $N$}}
	\put(7,110){\makebox[0pt][l]{\small $Y$}}
	\put(-7,113){\makebox[0pt][r]{\small pounds}}
	\put(-7,94){\makebox[0pt][r]{\small per tree}}
	\put(-6,67){\makebox[0pt][r]{\small 750}}
	\put(260,-18){\makebox[0pt]{\small trees}}
	\multiput(0,70)(5,0){8}{\circle*{1}}
	\multiput(40,70)(3.25,-1.625){43}{\circle*{1}}
	\multiput(40,70)(0,-13){5}{\circle*{1}}
	\put(40,-18){\makebox[0pt]{\small 40}}
	\put(180,-18){\makebox[0pt]{\small 180}}
	
\end{picture}
}
\end{picture}
\setlength{\unitlength}{1pt}
}


When there are only a few trees, they don't get in each other's way,
and they produce at the maximum level---say, 750 pounds per tree.
Hence the graph starts off level.  At some point, the trees become too
crowded to produce {\em anything\/}!  In between, the yield per tree
drops off as shown by the curved middle part of the graph.

We want to choose $N$ so that the {\em total yield}, $T$, will be as
large as possible.  We have $T(N) = Y(N) \cdot N$, but since we don't
know $Y(N)$ very precisely, it is difficult to analyze $T(N)$.  To
help, let's replace $Y(N)$ by the approximation shown in the graph on
the right.  Now carry out an analysis using this graph to represent
$Y(N)$.

\num Find a formula for the straight segment of the new graph of
$Y(N)$ on the interval $40 \leq N \leq 180$.  What is the formula for
$T(N)$ on the same interval?

\num What are the formulas for $T(N)$ when $0 \leq N \leq 40$ and when
$180 \leq N$?  Graph $T$ as a function of $N$.  Describe the graph in
words.

\num What is the maximum possible total yield $T$?  For which $N$ is
this maximum attained?

\num Suppose the endpoints of the sloping segment were $P$ and $Q$,
instead of 40 and 180, respectively.  Now what is the formula for
$T(N)$?  (Note that $P$ and $Q$ are {\em parameters\/} here.
Different values of $P$ and $Q$ will give different models for the
behavior of the total output.)  How many trees would then produce the
maximum total output?  Expect the maximum to depend on the parameters
$P$ and $Q$.


\subsubsection{Rate Equations}
\index{rate equation}

Do the following exercises by hand.  You may wish to check your
answers by using suitable modifications of the program SIR.

\num {\bf Radioactivity}.\index{radioactivity} From exercise~25 of
section~2 we know a sample of $R$ grams of radium decays into lead at
the rate
\[
R' = \dfrac{-1}{2337} R \quad \mbox{grams per year} .
\]
Using a step size of 10 years, estimate how much radium remains in a
0.072 gram sample after 40 years.

\num {\bf Poland and Afghanistan}.\index{Poland}\index{Afghanistan} If
$P$ and $A$ denote the populations of Poland and Afghanistan,
respectively, then their net per capita growth rates imply the
following equations:
\begin{align*}
P' &= .009 \, P \quad \mbox{persons per year;} \\
A' &= .0216 \, A \quad \mbox{persons per year.} 
\end{align*}
(See exercise~23 of section~2.)  In 1985, $P = 37.5$ million, $A = 15$
million.  Using a step size of 1 year, estimate $P$ and $A$ in 1990.

\num {\bf Falling bodies.}  If $d$ and $v$ denote the distance fallen
(in feet) and the velocity (in feet per second) of a falling body,
then the motion can be described by the following equations:
\begin{align*}
d' &= v  \quad \mbox{feet per second};\\
v' &= 32 \quad \mbox{feet per second per second}.
\end{align*}
(See exercise~21 of section~2.)  Assume that when $t=0$, $d=0$ feet and
$v=10$ feet/sec.  Using a step size of 1~second, estimate $d$ and $v$
after 3~seconds have passed.

\cleardoublepage

